% Appendix (unlimited length). Included from main.tex after \appendix\onecolumn.

\section*{Appendix overview}
\Cref{app:exact} gives full proofs of \cref{sec:exact}, \cref{app:enum} the
exhaustive enumeration check, \cref{app:agg} the proofs of \cref{sec:agg},
\cref{app:finite} the proofs of \cref{sec:finite}, \cref{app:sims} the
simulation checks of F1--F3, \cref{app:lean} the map from statements to Lean
names, \cref{app:real} details of the real-data studies, \cref{app:estimand}
the real-data estimand $\tilde s$, \cref{app:prereg}
the analysis-plan timeline, prior exposure and deviations, \cref{app:e1,app:e2} the
experiments E1 and E2 in full, \cref{app:noloop,app:lags,app:er,app:x3} the
loop-free graphs, lagged unrollings, the random summary graphs (with the full
proof of \cref{prop:er} in \cref{app:erproof}) and the accuracy interpretation.

Status labels. \emph{Machine-checked}: a Lean theorem with no \texttt{sorry}
and only the standard axioms. \emph{Proved}: a complete paper proof given
here. \emph{Asymptotic}: a limit statement proved here from the multivariate
Breuer--Major theorem \citep{arcones1994limit,nourdin2011quantitative} (cited, not
re-proved). \emph{Sketched}: a derivation with one analytic step left out,
stated as such (used only for statements outside the main claims: the windowed
loop-free stationary value, the layered-DAG stationary value and the constants
of the lagged recovery bound).
\emph{Enumerated}: checked by exact computation on many instances, not proved.

%=============================================================================
\section{Proofs for Section~\ref{sec:exact}}
\label{app:exact}

\textbf{Definition of $\srho$, verbatim.} \citet[Sec.~2.2]{reisach2026topological}
(arXiv v1) write: ``For a graph $G=(V,E)$ with $|E|\ge1$ we define sortability
by the number of relatives, or rel-sortability for short, as
$s_\rho(G):=\frac{\sum_{(x,y)\in E}s_\rho(x,y)}{|E|}$ with
$s_\rho(x,y):=0$ if $\rho(x)>\rho(y)$, $\frac12$ if $\rho(x)=\rho(y)$, $1$ if
$\rho(x)<\rho(y)$, for $(x,y)\in E$.'' (their Eq.~(4); the cases are typeset as a
brace there), and add: ``Note that our definition differs from the sortabilities
in these works in that we count edges instead of paths.'' Our $\srho$ is this
edge average; in $\GT$ the edges are the copies $(i(t),j(t{+}1))$.

Throughout, $\GS=(V,\Es)$ has a self-loop at every node, $p=|V|\ge2$, and
$d(u,v)$ is as in \cref{sec:setting}. We use two facts of
\citet{reisach2026topological}, proved there for DAGs: (Thm.~1) if
$x\in\anc(y)$ then $\rel(x)\subseteq\rel(y)$; (Lemma~1, eq.~(10)) if
$x\in\anc(y)$ then $\rel(x)\subsetneq\rel(y)$ iff $\mathrm{roots}(y)\setminus
\mathrm{roots}(x)\ne\emptyset$, where $\mathrm{roots}(y)=\{x\in\anc(y):\anc(x)=\{x\}\}$.
Both are also machine-checked independently
(\lname{relSet_subset_of_edge}, and the direct characterisation
\lname{isRel_unrolled_iff}).

\begin{lemma}[ancestors in $\GT$]\label{lem:anc}
$u(s)\in\anc(v(t))$ iff $s\le t$ and $d(u,v)\le t-s$.
\lean{\lname{unrolled_reach_iff}, \lname{unrolled_reach_iff_dist}.}
\end{lemma}
\begin{proof}
Every edge of $\GT$ increases time by one, so a directed path from $u(s)$ to
$v(t)$ has $t-s$ edges, and its projection to $V$ is a walk of length $t-s$
from $u$ to $v$ in $\GS$ (self-loops allowed as steps). Such a walk exists iff
$d(u,v)\le t-s$: a walk contains a path, so its length is at least
$d(u,v)$; conversely, pad a shortest path with $t-s-d(u,v)$ self-loop steps at
$u$. Every such walk lifts to a directed path in $\GT$.
\end{proof}

\begin{lemma}[roots]\label{lem:roots}
$\mathrm{roots}(v(t))=\{u(1):d(u,v)\le t-1\}$, and
$\desc(u(1))=\{w(s):d(u,w)\le s-1\}$.
\lean{\lname{roots_unrolled_iff}, \lname{isRoot_unrolled_iff}, \lname{desc_root_unrolled_iff}.}
\end{lemma}
\begin{proof}
A node $u(s)$ with $s\ge2$ has the parent $u(s-1)$ through the self-loop, so
it is not a root; $u(1)$ has no parent. Apply \cref{lem:anc} with $s=1$, in
both directions.
\end{proof}

\begin{proof}[Proof of \cref{thm:edge}]
$i(t)$ is a parent of $j(t{+}1)$, so $\rel(i(t))\subseteq\rel(j(t{+}1))$ and
$\srho(e)\in\{1/2,1\}$, with value $1$ iff the inclusion is strict, iff
$\mathrm{roots}(j(t{+}1))\setminus\mathrm{roots}(i(t))\neq\emptyset$. By
\cref{lem:roots} this set is $\{u(1):d(u,j)\le t,\ d(u,i)>t-1\}$, and since
$d$ is integer-valued or $\infty$, $d(u,i)>t-1$ iff $d(u,i)\ge t$.
\end{proof}

\begin{proof}[Proof of \cref{prop:boundary} (and the general count)]
Since every ancestor of $y$ descends from a root of $y$ that is itself an
ancestor, $\rel(y)=\desc(\mathrm{roots}(y))$. By \cref{lem:roots},
$w(s)\in\rel(v(t))$ iff some $u$ has $d(u,v)\le t-1$ and $d(u,w)\le s-1$,
i.e.\ iff $s\ge1+m_t(v,w)$ with $m_t(v,w)=\min\{d(u,w):d(u,v)\le t-1\}$.
Counting $s\in\{1,\dots,T\}$ gives
$\rho(v(t))=\sum_{w}\max(T-m_t(v,w),0)$ (\lname{rho_unrolled_eq}). For
$t\ge p$ every finite $d(u,v)$ is at most $p-1\le t-1$, so the admissible $u$
are $\anc_S(v)$ and $m_t(v,w)=d(\anc_S(v),w)$, finite iff $w\in\rel_S(v)$ and
then at most $p-1<t\le T$. (The Lean statement \lname{rho_unrolled_large_t}
needs only $t\ge p$; an earlier hypothesis $T\ge p$ is redundant.)
\end{proof}

\begin{proof}[Proof of \cref{thm:limit}]
Crossing edge $(i,j)\in\Ex$: $j$ does not reach $i$, so $d(j,i)=\infty$, and
$u=j$ gives $d(j,j)=0\le t\le\infty$ for every $t$: all $T-1$ copies score
$1$. Edge in $\Eo$: $i$ and $j$ reach each other (or $i=j$), so for every $u$,
$d(u,i)<\infty$ iff $d(u,j)<\infty$. If the condition of \cref{thm:edge} held
with $d(u,i)=\infty$ then $d(u,j)=\infty>t$, a contradiction; hence it forces
$t\le d(u,i)\le p-1$, and at most $N_T(i,j)\le\min(p-1,T-1)$ copies score
$1$. Summing over the $(T-1)|\Es|$ edges,
\[
\srho(\GT)=\tfrac12+\frac{1}{2(T-1)|\Es|}\Big[(T-1)|\Ex|+\sum_{(i,j)\in\Eo}N_T(i,j)\Big],
\]
which gives both bounds; let $T\to\infty$.
\end{proof}

\begin{proof}[Proof of \cref{cor:iff}]
$\Ex=\emptyset$ iff no edge joins two SCCs. A weakly connected component made of
two or more SCCs is connected by undirected edges, one of which joins two
different SCCs; conversely a strongly connected component has no crossing
edge. The uniform bound is \cref{thm:limit} with $|\Ex|\ge1$. For the DAG case,
every SCC is a singleton, $\Ex$ consists of the $m$ non-loop edges, and
$|\Es|=m+p$.
\end{proof}

\begin{proof}[Proof of \cref{thm:stat}]
In $G_{B+T}$, by \cref{lem:roots} shifted by $B$, $w(s)\in\rel(v(t))$ iff some
$u$ has $d(u,v)\le t-1$ and $d(u,w)\le s-1$. For window nodes $t-1,s-1\ge
B\ge p-1$, and every finite distance is at most $p-1$, so the condition is
$u\in\anc_S(v)\cap\anc_S(w)$, independent of $s,t$. Hence
$\rel_W(v(t))=\{w(s)\in W: w\in\rel_S(v)\}$, of size $T|\rel_S(v)|$. On
$V\times\mathbb Z$, \cref{lem:anc} holds verbatim and the common ancestor can
be taken arbitrarily early, giving the same set. Each copy of $(i,j)$ therefore
compares $|\rel_S(i)|$ with $|\rel_S(j)|$; as $i\in\anc_S(j)$,
$\rel_S(i)\subseteq\rel_S(j)$, so the score is $1$ iff
$\rel_S(i)\ne\rel_S(j)$, and $1/2$ for self-loops.
Source SCCs: $\rel_S(j)\setminus\rel_S(i)\ne\emptyset$ iff
$\anc_S(j)\not\subseteq\rel_S(i)$ (if $z\in\rel_S(i)$ then
$\desc_S(z)\subseteq\rel_S(i)$). Every node of $\anc_S(j)$ is reached from some
$\sigma\in\Src(j)$, and $\desc_S(\sigma)\subseteq\rel_S(i)$ whenever
$\sigma\subseteq\rel_S(i)$. A source SCC meets $\rel_S(i)$ iff it reaches $i$
(its only ancestors are its own nodes), and then it is contained in
$\rel_S(i)$. So the condition is $\Src(j)\not\subseteq\Src(i)$, which, as
$\Src(i)\subseteq\Src(j)$, is $\Src(i)\ne\Src(j)$.
\end{proof}

\begin{proof}[Proof of \cref{cor:statiff}]
If each weak component $C$ has a unique source SCC $\sigma_C$, then
$\Src(v)=\{\sigma_C\}$ on $C$ and $\Eu=\emptyset$. Conversely, if
$\Eu=\emptyset$, $\Src$ is constant along edges and hence on each weak
component; a node $v$ of a source SCC $\sigma$ has $\Src(v)=\{\sigma\}$, so
there is only one. Out-forests (at most one parent besides itself; acyclicity
is not needed, \lname{stat_half_of_outForest}) and strongly connected graphs
(\lname{stat_half_of_strong}) have one source SCC per component. In an
in-forest (acyclic apart from self-loops, out-degree at most one), if
$u\in\anc(v)$ then $\desc(u)$ is the unique downstream path through $v$, so
$\rel(v)=\anc(v)\cup\desc(v)$ (\lname{inForest_isRel_iff}); the source SCCs
are the headwaters, and $\Src(j)=\Src(i)\cup\bigcup_{k}\Src(k)$ over the other
parents $k$ of $j$, with disjoint headwater sets on different branches, so
$\Src(j)\ne\Src(i)$ iff $j$ is a confluence (\lname{inForest_up_iff}).
Out-degree at most one alone is not enough: for $a\rightleftarrows b$ with
$c\to a$, $a$ has two parents but the only source SCC is $\{c\}$.
\end{proof}

\begin{remark}[self-loops]
Without self-loops the lemmas fail: for the single edge $0\to1$ without loops,
$\srho(\GT)=1/2$ for every $T\ge2$ (\lname{edgeNoLoops_sRho}); with loops the
limit is $2/3$ (\lname{edgeLoops_limit}).
\end{remark}

%=============================================================================
\section{Exhaustive Enumeration Check}
\label{app:enum}

Before formalisation, every statement of \cref{sec:exact} was checked by
brute force (bitset relatives with the definitions of
\citealp{reisach2026topological}) on all $4+64+4096$ summary graphs with
$p=2,3,4$ (all $T\le9$) and on $10{,}000$ random graphs ($p=2$--$7$; Erd\H{o}s--R\'enyi,
DAG, in-forest, out-forest, chain with back edge and block-SCC families;
three values of $T\le14$ each, fixed seed). In total 46{,}724 (graph,~$T$)
cases and 1{,}895{,}609 edges were checked with zero failures of
\cref{thm:edge}, \cref{prop:boundary}, the bounds of \cref{thm:limit},
\cref{thm:stat} (with $B=p-1$ and $B=2p+1$) and the source-SCC
characterisation. The conjecture fixed in plan 1 (``an edge is strictly
sorted in the limit iff $|\rel_S(j)|>|\rel_S(i)|$'') disagreed with the exact
$\GT$ limit in 9{,}603 random graphs and never with the stationary value; the
conjecture is \cref{thm:stat}, not the limit of $\GT$.

%=============================================================================
\section{Proofs for Section~\ref{sec:agg}}
\label{app:agg}

Assumptions: $X(t)=\sum_{\tau=1}^{L}A_\tau X(t-\tau)+\varepsilon(t)$, $t\in\mathbb Z$,
i.i.d.\ innovations with diagonal positive-definite covariance $\Sigma$, stable
(all roots of $\det(I-\sum_\tau A_\tau z^\tau)$ outside the unit disc),
stationary solution; the summary graph has $i\to j$ iff $(A_\tau)_{ji}\ne0$ for
some $\tau$. $\Gamma(h)$ is the autocovariance, $A(1)=\sum_\tau A_\tau$, and for a
block size $D$ the block means are $\bar Y_s=\frac1D\sum_{u<D}X(sD{+}u)$ and the
subsamples $Z_s=X(sD)$.

\begin{proposition}[relatives survive any per-node filter]\label{prop:zero}
If $\anc_S(i)\cap\anc_S(j)=\emptyset$ then $X_i$ and $X_j$ are independent at
all lags, and so are any per-node, time-invariant (possibly nonlinear)
transforms of them, including $\bar Y$ and $Z$ for every $D$. If the
ancestor sets intersect, $\Gamma(0)_{ij}$, $\mathrm{Cov}(\bar Y)_{ij}$ and the
long-run covariance are nonzero for Lebesgue-almost every parameter.
\nolean{needs the SVAR moving-average representation and real-analytic genericity}
\end{proposition}

\begin{theorem}[block means vs subsampling]\label{thm:agg}
(a) $\|D\,\mathrm{Cov}(\bar Y)-S_{LR}\|\le\frac1D\sum_h|h|\|\Gamma(h)\|$,
with $S_{LR}=(I-A(1))^{-1}\Sigma(I-A(1))^{-\top}$.
(b) If $\GS$ is acyclic apart from self-loops, $S_{LR}$ is the covariance of a
static SEM (this long-run SEM is due to \citealp[Eqs.~(5), (10)]{gong2017aggregated}) on $\GS$ with weights $A(1)_{ji}/(1-\lambda_j)$ and noise variances
$\sigma_j^2/(1-\lambda_j)^2$, $\lambda_j=A(1)_{jj}<1$.
(c) $\mathrm{Cov}(Z_s)=\Gamma(0)$ for every $D$.
(d) Hence, for large $D$, block-mean $R^2$-sortability is $0$ on out-stars and
$1$ on in-stars with $\ge2$ arms, for all weights; under subsampling every
contemporaneous criterion keeps its $\Gamma(0)$ value.
\lean{(a) counting identity \lname{fejer_identity}; (b) \lname{longrun_factor}, \lname{selfloop_sum_lt_one}; the probabilistic steps, the rate and (c)--(d) are not machine-checked.}
\end{theorem}
Part (a) is the standard Fej\'er/Ces\`aro computation for the variance of the
mean of a stationary series \citep[Ch.~7]{brockwell1991time}, written for
block means of a vector process; (b) is due to \citet{gong2017aggregated}; (c)
is stationarity. We claim only the sortability consequences (d).

\begin{proof}[Proof of \cref{prop:zero}]
By the MA($\infty$) representation $X(t)=\sum_k\Psi_k\varepsilon(t-k)$, and
because $\Psi_k$ is supported on the reachability relation, $X_i(t)$ is a
measurable function of $\{\varepsilon_k(s):k\in\anc_S(i)\}$. Disjoint ancestor
sets mean independent innovation sets, so the processes $X_i$ and $X_j$ are
independent, and so are any per-node transforms. For the generic part, each
covariance entry is real-analytic in $(A,\Sigma)$ on the connected open set of
stable parameters with the given support, hence either identically zero or
nonzero almost everywhere. With all coefficients positive, including the
self-loop at a common ancestor $k$, all $\Psi$ entries are nonnegative and
$\Psi_{ik}(m),\Psi_{jk}(m)>0$ for all large $m$ (self-loops pad the walks), so
$\Gamma(0)_{ij}>0$; the same holds for $\Gamma(h)$, $h\ge0$, and therefore for
the Fej\'er sums and $S_{LR}$. ``Generic'' does not mean ``large'': with little
autocorrelation, $\Gamma(0)_{ij}$ along a lagged edge can be tiny (down to
$10^{-14}$ in our exact computations; exactly $0$ without autocorrelation).
\end{proof}

\begin{proof}[Proof of \cref{thm:agg}]
(a) $D^2\mathrm{Cov}(\bar Y)=\sum_{u,w<D}\Gamma(u-w)=D\Gamma(0)+\sum_{h=1}^{D-1}(D-h)(\Gamma(h)+\Gamma(-h))$
(the counting identity is \lname{fejer_identity}). The stable VAR has
$\|\Gamma(h)\|\le C r^{|h|}$, $r<1$, so $\sum_h|h|\|\Gamma(h)\|<\infty$ and
$S_{LR}=\sum_h\Gamma(h)=\Psi(1)\Sigma\Psi(1)^\top$ with
$\Psi(1)=(I-A(1))^{-1}$. Then
$S_{LR}-D\,\mathrm{Cov}(\bar Y)=\sum_{|h|\ge D}\Gamma(h)+\sum_{|h|<D}\frac{|h|}{D}\Gamma(h)$,
and $1\le|h|/D$ in the first sum.
(b) Write $I-A(1)=(I-\Lambda)-N$ with $\Lambda$ the diagonal and $N$ the
off-diagonal part; then $I-A(1)=(I-\Lambda)(I-\tilde B)$ with
$\tilde B=(I-\Lambda)^{-1}N$ strictly triangular in a topological order
(\lname{longrun_factor}, \lname{longrun_inv}). Under acyclicity the AR
polynomial factorises into univariate $\phi_j(z)=1-\sum_\tau(A_\tau)_{jj}z^\tau$;
stability, $\phi_j(0)=1$ and the intermediate value theorem give
$\phi_j(1)=1-\lambda_j>0$ (\lname{selfloop_sum_lt_one}).
(c) Stationarity.
(d) For a linear SEM $X=BX+e$ with noise variances $\omega$,
$R^2_i=1-1/(\Sigma_{ii}(\Sigma^{-1})_{ii})$ and
$(\Sigma^{-1})_{ii}=1/\omega_i+\sum_{k\in\mathrm{ch}(i)}B_{ki}^2/\omega_k$. In
an out-star with root $0$ and leaves $k$, $1-R^2_0=1/(1+\sum_kr_k)$ and
$1-R^2_k=1/(1+r_k)$ with $r_k=\omega_0B_{k0}^2/\omega_k>0$, so with $\ge2$
leaves the root has strictly larger $R^2$ than every leaf; the in-star is
symmetric. Strict inequalities in the limit SEM of (b), plus (a) and
continuity of $S\mapsto R^2$ at positive-definite $S$, give the claim for all
large $D$. Under subsampling every contemporaneous statistic is a function of
$\Gamma(0)$.
\end{proof}

The $R^2$ statement for stars strengthens the three-node analysis in
\citet[App.~D]{herman2025unitless}, whose qualitative message (that $R^2$
reflects connectivity, with an asymmetry between colliders and confounders)
we credit. Numerical observations (exact population moments, not proved):
mean block-mean $R^2$-sortability over 150 in-trees moves from $0.853$
($D=1$) to $0.834$ ($D=4$--$8$) and $0.880$ ($D=256$), so it is not monotone in
$D$; on out-stars subsampled data are $R^2$-sorted at $0.842$ on average while
block means give $0.546$ ($D=4$), $0.101$ ($D=16$) and $0$ in the limit.

%=============================================================================
\section{Proofs for Section~\ref{sec:finite}}
\label{app:finite}

Notation. $\hat C_{ab}=\frac1n\sum_t(X_a(t)-\bar X_a)(X_b(t)-\bar X_b)$,
$\hat r_{ab}=\hat C_{ab}/\sqrt{\hat C_{aa}\hat C_{bb}}$. Constants:
$B_{ab}$ as in the main text; $D_\gamma=\max_a2\sum_h\gamma_{aa}(h)^2$;
$\kappa_1=\max_a\sum_h|\gamma_{aa}(h)|$;
$\kappa_2=\max_{a\ne b}(\sum_h|\gamma_{aa}(h)|+\sum_h|\gamma_{ab}(h)|)$. For
i.i.d.\ sampling, $B_{ab}=1+\rho_{ab}^2$, $D_\gamma=2$, $\kappa_1=1$.

\begin{lemma}[Gaussian chaos tail]\label{lem:chaos}
Let $g\sim N(0,I_N)$, $A$ symmetric, $S=g^\top Ag-\mathrm{tr}A$. For $x>0$,
$P(S\ge2\|A\|_F\sqrt x+2\|A\|_{op}x)\le e^{-x}$, and the same for $-S$.
\lean{algebraic steps \lname{phi_le}, \lname{chernoff_logmgf_le}, \lname{chernoff_choice}, \lname{sqrt_two_v_x}; the probability statement is not machine-checked.}
\end{lemma}
\begin{proof}
Diagonalise; $S$ has the law of $\sum_i\lambda_i(g_i^2-1)$. For $2s\lambda<1$,
$\log Ee^{s\lambda(g^2-1)}=\varphi(s\lambda)$ with
$\varphi(u)=-\frac12\log(1-2u)-u=\sum_{k\ge2}(2u)^k/(2k)\le u^2/(1-2|u|)$ for
$|u|<1/2$. With $c=2\|A\|_{op}$, $v=2\|A\|_F^2$ and $0\le s<1/c$,
$\log Ee^{sS}\le vs^2/(2(1-cs))$. Markov's inequality with $t=\sqrt{2vx}+cx$,
$q=\sqrt{2x/v}$ and $s=q/(1+cq)$ gives exponent exactly $x$. This is the
Gaussian case of \citet[Prop.~1]{hsu2012tail}, who derive it from
\citet{laurent2000adaptive}; the proof here is self-contained.
\end{proof}

\begin{lemma}[exact variance]\label{lem:var}
With $\tilde\gamma_{ab}=\frac1n\sum_tX_a(t)X_b(t)=g^\top Ag$ for
$A=C^{1/2}MC^{1/2}$, $M=\frac1{2n}\begin{psmallmatrix}0&I\\I&0\end{psmallmatrix}$,
\[
2n\|A\|_F^2=\sum_{|h|<n}\Big(1-\tfrac{|h|}{n}\Big)\big[\gamma_{aa}(h)\gamma_{bb}(h)+\gamma_{ab}(h)\gamma_{ab}(-h)\big]\le B_{ab},
\]
$\|A\|_{op}\le\kappa_2/(2n)$; on the diagonal $2n\|A\|_F^2\le D_\gamma$ and
$\|A\|_{op}\le\kappa_1/n$; and $\bar X_a\sim N(0,\bar\sigma_a^2)$ with
$\bar\sigma_a^2\le\kappa_1/n$.
\lean{trace and Toeplitz algebra \lname{quad_M}, \lname{trace_MC}, \lname{lemma12_trace}, \lname{lemma12_fejer}, \lname{lemma12_le_B}, \lname{lemma12_diag_le_D}; the step $\|C^{1/2}MC^{1/2}\|_F^2=\mathrm{tr}(MCMC)$ and the operator-norm bounds are not machine-checked.}
\end{lemma}
\begin{proof}
$\|A\|_F^2=\mathrm{tr}(MCMC)$ and, with Toeplitz blocks
$(C_{ab})_{st}=\gamma_{ab}(s-t)$,
$4n^2\mathrm{tr}(MCMC)=2\sum_{s,t}[\gamma_{aa}(s-t)\gamma_{bb}(s-t)+\gamma_{ab}(s-t)\gamma_{ab}(t-s)]$
(using that $\gamma_{aa}$, $\gamma_{bb}$ are even); count pairs with $s-t=h$.
$\|A\|_{op}\le\|C\|_{op}\|M\|_{op}$ and $\|C\|_{op}$ is at most the largest
absolute row sum.
\end{proof}

\begin{lemma}[deviation event]\label{lem:dev}
Let $\varepsilon_{ab}=\sqrt{2B_{ab}x/n}+\kappa_2x/n+2\kappa_1y/n$ and
$\varepsilon_d=\sqrt{2D_\gamma y/n}+4\kappa_1y/n$. With probability at least
$1-d(d{-}1)e^{-x}-4de^{-y}$, $|\hat C_{ab}-\rho_{ab}|\le\varepsilon_{ab}$ for all
$a\ne b$ and $|\hat C_{aa}-1|\le\varepsilon_d$ for all $a$.
\lean{deterministic steps \lname{centred_cov_eq}, \lname{centred_dev_offdiag}, \lname{centred_dev_diag}, \lname{lemma13_eps}, \lname{lemma13_count}.}
\end{lemma}
\begin{proof}
$\hat C_{ab}=\tilde\gamma_{ab}-\bar X_a\bar X_b$. Apply \cref{lem:chaos,lem:var}
to each off-diagonal and diagonal $\tilde\gamma$, bound $|\bar X_a|\le
\sqrt{2\kappa_1y/n}$ by the Gaussian Chernoff bound, and take a union bound.
\end{proof}

\begin{lemma}[correlation perturbation]\label{lem:pert}
If $a',b'\in[1-\varepsilon_d,1+\varepsilon_d]$, $\varepsilon_d<1$ and
$|c'-\rho|\le\varepsilon$, then
$|c'/\sqrt{a'b'}-\rho|\le(\varepsilon+|\rho|\varepsilon_d)/(1-\varepsilon_d)$.
\lean{\lname{corr_perturb} (the hypothesis $|\rho|\le1$ of our earlier statement is not needed).}
\end{lemma}
\begin{proof}
$|c'/\sqrt{a'b'}-\rho|\le|c'-\rho|/\sqrt{a'b'}+|\rho||1/\sqrt{a'b'}-1|$, with
$1/\sqrt{a'b'}\le1/(1-\varepsilon_d)$ and $|1/\sqrt{a'b'}-1|\le\varepsilon_d/(1-\varepsilon_d)$.
\end{proof}

\begin{theorem}[recovery bound]\label{thm:f1}
For $x,y>0$ let $\varepsilon_{ab}$, $\varepsilon_d<1$ be as in \cref{lem:dev} and
$\delta_{ab}=(\varepsilon_{ab}+|\rho_{ab}|\varepsilon_d)/(1-\varepsilon_d)$. If
thresholds separate ($\delta_{ab}\le\tau_{ab}$ on non-relatives and
$|\rho_{ab}|-\delta_{ab}>\tau_{ab}$ on relatives), the declared set equals the true
one, and hence $\hat s=\srho^{\mathrm{pop}}$, with probability at least
$1-d(d{-}1)e^{-x}-4de^{-y}$. For a common threshold $\tau<\rho_{\min}$ the
failure probability decays as
$\exp(-n\min(\tau,\rho_{\min}-\tau)^2/(2(\sqrt B+\rho_{\max}\sqrt{D_\gamma})^2))$.
\lean{deterministic core \lname{thm15_det}, \lname{decide_eq_target}, \lname{corr_perturb}, \lname{sOf_cntDecl_congr}; Chernoff and trace algebra \lname{chernoff_choice}, \lname{phi_le}, \lname{lemma12_trace}, \lname{lemma12_le_B}. The Gaussian tail bound itself is not machine-checked.}
\end{theorem}
The proof combines a Gaussian-chaos tail bound (the Gaussian case of
\citealp[Prop.~1]{hsu2012tail}; see also \citealp{laurent2000adaptive}), an exact
finite-$n$ variance identity for $\hat r$, and a deterministic perturbation
lemma.

\begin{proof}[Proof of \cref{thm:f1}]
Set $\delta_{ab}=(\varepsilon_{ab}+|\rho_{ab}|\varepsilon_d)/(1-\varepsilon_d)$.
On the event of \cref{lem:dev}, \cref{lem:pert} gives
$|\hat r_{ab}-\rho_{ab}|\le\delta_{ab}$, so non-relatives satisfy
$|\hat r|\le\tau$ and relatives $|\hat r|>\tau$ (\lname{decide_eq_target},
\lname{thm15_det}). Equal declared sets give equal counts and hence
$\hat s=\srho^{\mathrm{pop}}$ (\lname{sOf_cntDecl_congr}). For a VAR observed
as a long stationary series, $\srho^{\mathrm{pop}}$ is the stationary value
of \cref{thm:stat} on non-self-loop edges (\cref{prop:zero}). Explicit form:
with $y=x+\log(4/(d-1))_+$ (\lname{explicit_y}) and a common threshold
$\tau<\rho_{\min}$,
$P(\hat R\ne R)\le2d(d{-}1)\exp\!\big(-n\min(\tau,\rho_{\min}-\tau)^2/(2(\sqrt B+\rho_{\max}\sqrt{D_\gamma})^2)(1+o(1))\big)$.
\end{proof}

\textbf{The multivariate Breuer--Major theorem.} We use it in the form stated
by \citet[Thm.~1.1]{nourdin2011quantitative}, who attribute the vector case to
\citet{arcones1994limit}: let $X=(X_k)_{k\in\mathbb Z}$ be a centred stationary
Gaussian process in $\mathbb R^d$ with covariances $r^{(i,l)}(j)=E[X^{(i)}_1X^{(l)}_{1+j}]$,
and $f:\mathbb R^d\to\mathbb R$ with $E f(X_1)^2<\infty$ and Hermite rank $q\ge1$
with respect to $X$ (orthogonal to all polynomials of degree $<q$, not to some of
degree $q$). If $\sum_{j\in\mathbb Z}|r^{(i,l)}(j)|^q<\infty$ for all $i,l$, then
$n^{-1/2}\sum_{k\le n}(f(X_k)-Ef(X_k))\Rightarrow N(0,\sigma^2)$ with
$\sigma^2=\mathrm{Var}f(X_1)+2\sum_{k\ge1}\mathrm{Cov}(f(X_1),f(X_{1+k}))\in[0,\infty)$.

\begin{proof}[Proof of \cref{prop:nominal}]
Take $d=2$, $X_k=(X_a(k),X_b(k))$ and, for $c\in\mathbb R^3$,
$f_c(x)=c_1(x_a^2-1)+c_2(x_b^2-1)+c_3(x_ax_b-\rho_{ab})$. $f_c$ is a centred
polynomial of degree $2$ without linear part, so it is orthogonal to every
polynomial of degree $\le1$ (odd Gaussian moments vanish); if it were orthogonal
to every polynomial of degree $2$ it would be orthogonal to itself, so either
$f_c(X_1)=0$ a.s.\ (then the partial sums vanish) or its Hermite rank is $2$.
The condition $\sum_h|\gamma_{ab}(h)|^2<\infty$ holds because the covariances are
absolutely summable and bounded. By the theorem and the Cram\'er--Wold device,
$\sqrt n(\tilde\gamma_{aa}-1,\tilde\gamma_{bb}-1,\tilde\gamma_{ab}-\rho_{ab})\Rightarrow N(0,V)$,
where, by Isserlis' theorem, $V(ab,cd)=\sum_h[\gamma_{ac}(h)\gamma_{bd}(h)+\gamma_{ad}(h)\gamma_{bc}(h)]$
(as in \cref{lem:var}). Centring by the sample means changes each entry by
$\bar X_a\bar X_b$, and $\sqrt n\bar X_a\bar X_b\to0$ in probability since
$E(\sqrt n\bar X_a)^2\le\kappa_1$ and $\bar X_b\to0$; Slutsky's lemma applies.
The delta method with gradient $(-\rho/2,-\rho/2,1)$ gives $\sqrt n(\hat r-\rho)
\Rightarrow N(0,W)$,
$W=V(ab,ab)-\rho[V(ab,aa)+V(ab,bb)]+\frac{\rho^2}{4}[V(aa,aa)+2V(aa,bb)+V(bb,bb)]$,
which equals $B_{ab}$ for an independent pair. Then
$\sqrt{n}\,\atanh\hat r\Rightarrow N(0,B_{ab})$ under the null and
$P(|\atanh\hat r|>z_c/\sqrt{n-3})\to2\bar\Phi(z_c/\sqrt{B_{ab}})$.
\end{proof}

\begin{proof}[Proof of \cref{prop:f2}]
$\mathrm{Cov}(X_u,X_v)=\gamma^Y_{uv}(0)+b_ub_v$ and $\gamma^Y_{uv}=0$ for
non-relatives (\cref{prop:zero}); divide by the standard deviations
(\lname{prop21_corr_std}). For relatives the sum is nonzero except on a
Lebesgue-null set. All counts equal $d$, so every edge scores $1/2$
(\lname{sOf_saturated}). The rate: apply \cref{thm:f1} with target ``all
pairs'' and $x=y=cn$, $c$ small enough that $\delta\le\kappa_{\min}/3$ and
$\tau_n\le\kappa_{\min}/3$ eventually; $|\hat s-1/2|\le1/2$ always
(\lname{abs_sOf_sub_half_le}). The onset: by the asymptotic normality of
$\atanh\hat r$, the declaration probability crosses $1/2$ where
$\tanh(z_c/\sqrt{n-3})=|\kappa|$, i.e.\ at $\nhalf$ (\lname{ncrit_eq_iff}).
\end{proof}

\begin{proposition}[conditioning]\label{prop:f3}
Let $Z=\sqrt r\,C+\sqrt{1-r}\,e$ with reliability $r=\mathrm{corr}(Z,C)^2$,
$\pi_{uv}$ the partial correlation given $Z$, and $\rho^Y$ the driver-free
correlation. Then
\[
\pi_{uv}=\frac{\rho^Y_{uv}\sqrt{(1{-}\ell_u)(1{-}\ell_v)}+(1{-}r)\,\mathrm{sgn}(b_ub_v)\sqrt{\ell_u\ell_v}}{\sqrt{(1-r\ell_u)(1-r\ell_v)}} .
\]
At $r=1$ the conditional relatives are the summary relatives and
\cref{thm:stat} is restored. For non-relatives with $\ell_u,\ell_v\in(0,1)$ and
$b_ub_v\ne0$, $|\pi_{uv}|$ is strictly decreasing in $r$ and vanishes only at
$r=1$; fewer than half of the non-relative pairs are declared only if
$r>r_{\mathrm{crit}}=1-\frac{\tau'(1-\ell)}{\ell(1-\tau')}$,
$\tau'=\tanh(z_c/\sqrt{n'-4})$ (equal shares). If the driver reaches nodes
through different filters, a scalar proxy leaves a nonzero residual even at
$r=1$. The reliability is identified by $\rho_{uZ}\rho_{vZ}/\rho_{uv}=r$, the
identity behind the three-indicator identification rule for one-factor models
\citep[Ch.~7]{bollen1989structural}.
\lean{\lname{prop31_std}, \lname{cor32}, \lname{absPi_strictAntiOn}, \lname{absPi_eq_zero_iff}, \lname{rcrit_iff}, \lname{thm35_bracket}, \lname{pc_lipschitz}, \lname{prop36_ar1_ne_zero}, \lname{prop36_rank_one}, \lname{prop38_ratio}.}
\end{proposition}

\begin{proof}[Proof of \cref{prop:f3}]
$\mathrm{Cov}(X_u,Z)=b_u\sqrt r$ and $\mathrm{Var}Z=1$; the partial covariance
is $\gamma^Y_{uv}+(1-r)b_ub_v$ and the partial variance
$\mathrm{Var}Y_u+(1-r)b_u^2$; divide and standardise (\lname{prop31},
\lname{prop31_std}; positive $\mathrm{Var}Y_u$ is required). At $r=1$,
$\pi=\rho^Y$ (\lname{cor32}). For non-relatives
$|\pi_{uv}|=\sqrt{\ell_u\ell_v}\sqrt{f_uf_v}$ with $f_w(r)=(1-r)/(1-r\ell_w)$
and $f_w'(r)=(\ell_w-1)/(1-r\ell_w)^2<0$ (\lname{absPi_strictAntiOn},
\lname{absPi_eq_zero_iff}; both need $b_ub_v\ne0$). Bounds
$(1-r)\sqrt{\ell_u\ell_v}\le|\pi|\le(1-r)\sqrt{\ell_u\ell_v}/(1-\ell_{\max})$
(\lname{absPi_lower}, \lname{absPi_upper}). $r_{\mathrm{crit}}$ solves
$(1-r)\ell/(1-r\ell)=\tau'$ (\lname{rcrit_iff}, \lname{rcrit_example}). The
finite-sample bracket uses the Lipschitz bound
$|g(a',b',c')-g(a,b,c)|\le L(q)\max(|a'-a|,|b'-b|,|c'-c|)$ for
$g(a,b,c)=(a-bc)/\sqrt{(1-b^2)(1-c^2)}$, $|b|,|c|\le q<1$,
$L(q)=1/(1-q^2)+4q/(1-q^2)^2$ (\lname{pc_lipschitz}, \lname{thm35_bracket}),
combined with \cref{lem:dev} for $(X,Z)$. Multi-filter driver: for
$D_v(t)=b_vC(t-k_v)$ with $C$ an AR(1) with coefficient $\phi$ and $Z=C(t)$,
\[
\pi_{uv}=\frac{b_ub_v(\phi^{|k_u-k_v|}-\phi^{k_u+k_v})}{\sqrt{(s_u^2-b_u^2\phi^{2k_u})(s_v^2-b_v^2\phi^{2k_v})}},
\]
nonzero for $k_u,k_v>0$, $\phi\in(0,1)$, $b_ub_v\ne0$ (\lname{prop36_ar1},
\lname{prop36_ar1_ne_zero}); in general $\pi_{uv}\propto\mathrm{Cov}(D_u-P_ZD_u,D_v-P_ZD_v)$
vanishes for all non-relative pairs iff the residual driver components are
uncorrelated, which holds for a rank-one driver with $Z=F$
(\lname{prop36_pcov}, \lname{prop36_zero_iff}, \lname{prop36_rank_one}). A
proxy that contains summary relatives of $u,v$ adds the terms
$-c_Z(b_uq_v+b_vq_u)-q_uq_v$ to the numerator (\lname{prop37_pcov},
\lname{prop37_bound}). Triad identity: $\rho_{uZ}\rho_{vZ}/\rho_{uv}=r$ and, for
$r>0$, $\ell_u=\rho_{uZ}^2/r$ (\lname{prop38_ratio}, \lname{prop38_loading}).
\end{proof}

%=============================================================================
\section{Simulation Checks of F1--F3}
\label{app:sims}

All with fixed seeds; full tables are in the supplementary material.
\textbf{F1} (5 processes: in-forests with $d=5$ and $d=10$, self-loop
persistence $\lambda\in\{0.5,0.9,0.97\}$, and a near-white-noise case;
$n=200$--$20{,}000$; 1{,}000 replications per cell; nominal, oracle-effective
and fixed thresholds). The bound was never violated in 105 cells. It is loose:
with a fixed threshold $\rho_{\min}/2$, it reaches $0.9$ at about 3--13 times
the $n$ at which the simulated probability does, and much later with
Bonferroni thresholds, where the Chernoff step loses the Mills-ratio factor.
The empirical $n\,\mathrm{Var}(\hat r)/W$ lies in $[0.93,1.11]$. With the
nominal test, $P(\hat R=R)$ plateaus at $0.94$, $0.22$, $0.03$ and $0.27$, as
\cref{prop:nominal} predicts (predicted vs simulated false-positive rates,
e.g.\ $0.33$ vs $0.32$--$0.33$ at $\lambda=0.9$).
\textbf{F2} ($\ell\in\{0.02,0.1,0.3,0.64\}$): the asymptotic detection curve
matches the simulated fraction of declared non-relatives within $0.02$ for
$n\ge200$; predicted $\nhalf=1{,}079$ for $\ell=0.1$, simulated $0.479$ at
$n=1{,}000$.
\textbf{F3}: predicted $n_{\mathrm{crit}}=472$ ($\ell=0.64$, $r=0.9$) gives a
simulated declared fraction $0.495$ at $n=471$; over 38 cells the asymptotic
prediction is within $0.03$. With travel-time lags $(0,1,2,4,8)$, $\phi=0.9$,
$r=1$, the declared fraction still rises to $0.63$ at $n=10{,}000$, as the
multi-filter part of \cref{prop:f3} predicts.

%=============================================================================
\section{Lean Formalisation: Map and Findings}
\label{app:lean}

Lean 4.35.0 and Mathlib v4.35.0; 3{,}963 lines in 16 files. The build log
(clean rebuild plus the incremental rebuild that added \lname{Observed.lean}),
the \texttt{\#print axioms} output for each of the 148 declarations in the axiom
report (all \texttt{[propext, Classical.choice, Quot.sound]}) and a grep showing
no \texttt{sorry}/\texttt{admit}/\texttt{axiom} are in the supplementary
material.

\textbf{What is and is not certified.} Certified: the exact and stationary
results of \cref{sec:exact}, the loop-free exact edge value, the accuracy
reading, the algebra of the estimand $\tilde s$ (\cref{app:estimand}), and the
deterministic and algebraic parts of \cref{sec:agg,sec:finite}. Not certified:
the ER limits (only the deterministic witness lemmas), the loop-free limit,
periodicity and stationary value, the lag-weighted extension, and every
probability bound, limit theorem and genericity statement.

\textbf{Breakdown of the 148 declarations} (counted from the Lean files):
29 headline theorems, i.e.\ Lean counterparts of numbered paper statements
(25 for \cref{sec:exact}, the loop-free edge value and the accuracy reading; 4
for \cref{prop:estimand}); 43 supporting graph lemmas (ancestry, roots, distance,
source SCCs, 5 ER witness lemmas, 6 estimand lemmas); 65 algebraic identities
and inequalities of the aggregation and finite-sample files (correlation
perturbation, Chernoff algebra, trace and Toeplitz identities, partial
correlation formulas, threshold locations; several are elementary, e.g.\
$|\hat s-\tfrac12|\le\tfrac12$); 11 examples (10 concrete graphs, 1 numerical
$r_{\mathrm{crit}}$ instance).

\begin{center}\small
\begin{tabular}{p{0.26\textwidth}p{0.52\textwidth}p{0.12\textwidth}}
\toprule
statement & Lean names & status \\
\midrule
\cref{lem:anc,lem:roots} & \lname{unrolled_reach_iff}, \lname{unrolled_reach_iff_dist}, \lname{roots_unrolled_iff}, \lname{desc_root_unrolled_iff}, \lname{isRel_unrolled_iff} & checked \\
\cref{thm:edge} & \lname{edge_score}, \lname{edge_score_dist}, \lname{strictCond_iff_dist} & checked \\
\cref{prop:boundary} & \lname{rho_unrolled_eq}, \lname{mT_le_iff}, \lname{rho_unrolled_large_t} & checked \\
\cref{thm:limit} & \lname{sRho_unrolled_exact}, \lname{sRho_unrolled_bounds}, \lname{tendsto_sRho_unrolled}, \lname{tendsto_sRho_GT} & checked \\
\cref{cor:iff} & \lname{tendsto_half_iff}, \lname{crossEdges_eq_empty_tfae}, \lname{sRho_unrolled_ge_of_cross}, \lname{reisach_theorem4}, \lname{limitValue_dag} & checked \\
\cref{thm:stat} & \lname{rhoBurn_eq}, \lname{sBurn_eq}, \lname{rhoZ_eq}, \lname{sZ_eq}, \lname{sOf_window}, \lname{relSet_eq_iff_src}, \lname{mem_upEdges_iff} & checked \\
\cref{cor:statiff} & \lname{sBurn_eq_half_iff}, \lname{sBurn_eq_half_iff_common_anc}, \lname{sZ_eq_half_iff}, \lname{upEdges_eq_empty_tfae}, \lname{stat_half_of_strong}, \lname{stat_half_of_outForest}, \lname{inForest_isRel_iff}, \lname{inForest_up_iff}, \lname{statValue_inForest}, \lname{statValue_inForest_gt_iff} & checked \\
\cref{tab:examples} & \lname{collider_limit}, \lname{fork_limit}, \lname{chain3_limit}, \lname{collider_stat}, \lname{fork_stat}, \lname{chain3_stat} & checked \\
self-loop remark & \lname{edgeLoops_limit}, \lname{edgeNoLoops_sRho} & checked \\
\cref{thm:noloop} (ancestors, relatives, edge value) & \lname{isRel_unrolled_iff_NL}, \lname{reduce_ancestor}, \lname{relSet_unrolled_eq_iff_NL}, \lname{edge_score_NL} & checked \\
\cref{cor:noloop} & \lname{sRho_unrolled_eq_NL}, \lname{sRho_unrolled_ge_half_NL} & checked \\
\cref{prop:noloopcross} (witness $u=j$, fresh parent) & \lname{strictNL_of_closed_walk}, \lname{strictNL_of_fresh_parent}, \lname{sRho_unrolled_eq_one_NL} & checked \\
loop-free examples (\cref{tab:noloop}) & \lname{colliderNL_sRho}, \lname{edgeLoopHead_sRho}, \lname{edgeNoLoops_sRho} & checked \\
eventual periodicity, limit, same-SCC edges, \cref{thm:noloopstat} & -- & paper / enumerated \\
\bottomrule
\end{tabular}
\end{center}

\begin{center}\small
\begin{tabular}{p{0.26\textwidth}p{0.52\textwidth}p{0.12\textwidth}}
\toprule
statement & Lean names & status \\
\midrule
\cref{prop:er} & witness lemmas \lname{upEdge_witness}, \lname{mem_upEdges_of_witness}, \lname{mem_upEdges_of_disjoint_anc}, \lname{upEdges_subset_confluenceNL}, \lname{card_upEdges_le_confluenceNL} & paper proof (\cref{app:erproof}); lemmas checked \\
\cref{prop:lags} & -- & paper only \\
\cref{prop:estimand} (a)--(d) & \lname{relObs_subset_of_edge}, \lname{score_obsEdge}, \lname{sTilde_eq}, \lname{sTilde_ge_half}, \lname{upObs_subset_upEdges}, \lname{sTilde_gt_half_iff}, \lname{exists_upEdge_of_sTilde_gt_half}, \lname{sTilde_eq_half_of_no_up}, \lname{mem_upObs_of_witness}, \lname{upObs_univ} & checked \\
\cref{prop:x3} (oracle sort accuracy; optimality) & \lname{accRule_sortRule}, \lname{accRule_sortRule_rho}, \lname{accRule_le_sOf}, \lname{accRule_le_sRho}, \lname{edgeAcc_sortRule_strict}, \lname{rho_not_gt_of_edge} & checked \\
\cref{prop:zero} & -- & paper only \\
\cref{thm:agg}(a) & \lname{fejer_identity} (counting identity) & partial \\
\cref{thm:agg}(b) & \lname{longrun_factor}, \lname{longrun_inv}, \lname{selfloop_sum_lt_one} & partial \\
\cref{thm:agg}(c,d) & -- & paper only \\
\cref{lem:chaos} & \lname{phi_le}, \lname{chernoff_logmgf_le}, \lname{chernoff_choice}, \lname{sqrt_two_v_x} & algebra only \\
\cref{lem:var} & \lname{quad_M}, \lname{trace_MC}, \lname{lemma12_trace}, \lname{lemma12_fejer}, \lname{avg_toeplitz_le_tsum}, \lname{lemma12_le_B}, \lname{lemma12_diag}, \lname{lemma12_diag_le_D} & algebra only \\
\cref{lem:dev} & \lname{centred_cov_eq}, \lname{centred_dev_offdiag}, \lname{centred_dev_diag}, \lname{lemma13_eps}, \lname{lemma13_count} & deterministic part \\
\cref{lem:pert} & \lname{corr_perturb} & checked \\
\cref{thm:f1} & \lname{decide_eq_target}, \lname{decide_eq_support}, \lname{thm15_det}, \lname{explicit_y}, \lname{sOf_cntDecl_congr} & deterministic part \\
\cref{prop:nominal} & -- & asymptotic (cited CLT) \\
\cref{prop:f2} & \lname{prop21_corr}, \lname{prop21_corr_std}, \lname{sOf_saturated}, \lname{sOf_const}, \lname{abs_sOf_sub_half_le}, \lname{thm22_const}, \lname{ncrit_iff}, \lname{ncrit_eq_iff} & deterministic part \\
\cref{prop:f3} & \lname{pcorr_eq_pcov}, \lname{pcorr_eq_corr_resid}, \lname{prop31}, \lname{prop31_std}, \lname{cor32}, \lname{prop33}, \lname{prop33_abs}, \lname{absPi_lower}, \lname{absPi_upper}, \lname{absPi_le_kappa}, \lname{absPi_lt_kappa}, \lname{absPi_eq_zero_iff}, \lname{absPi_strictAntiOn}, \lname{absPi_equal}, \lname{rcrit_iff}, \lname{rcrit_eq_iff}, \lname{rcrit_example}, \lname{pc_lipschitz}, \lname{thm35_bracket}, \lname{prop36_pcov}, \lname{prop36_zero_iff}, \lname{prop36_rank_one}, \lname{prop36_ar1_pcov}, \lname{prop36_ar1_pvar}, \lname{prop36_ar1}, \lname{prop36_ar1_ne_zero}, \lname{prop37_pcov}, \lname{prop37_bound}, \lname{prop38_ratio}, \lname{prop38_loading} & population part checked \\
\bottomrule
\end{tabular}
\end{center}

\textbf{Modelling choices.} Summary graphs are relations on a finite type with
a self-loop hypothesis; $\anc$ and $\desc$ are reflexive-transitive closures
(``a path exists iff a walk exists'' is the one convention adopted without
proof); $\GT$ with $T=n+1$ slices lives on $V\times\mathrm{Fin}(n+1)$; source
SCCs are represented by their member nodes (equality of SCC sets as equality of
node unions is elementary and not formalised). In the finite-sample files,
zero-mean variables are vectors of an arbitrary real inner-product space and
lag-0 covariance is the inner product; independence enters only as vanishing
inner products.

\textbf{Findings.} No statement was false. Formalisation added four missing
nondegeneracy hypotheses to our earlier finite-sample statements: (i) the
strict claims of \cref{prop:f3} need $b_ub_v\ne0$; (ii) $\ell_u=\rho_{uZ}^2/r$
needs $r>0$; (iii) nonvanishing of the multi-filter residual needs
$b_ub_v\ne0$; (iv) the partial correlation needs positive partial variances
($\mathrm{Var}Y_u>0$ suffices). It also showed that one hypothesis of
\cref{prop:boundary} and one of \cref{lem:pert} are redundant, and that the
in-forest statement needs acyclicity, not only out-degree at most one.

%=============================================================================
\section{Real-Data Details}
\label{app:real}

\textbf{CausalRivers} \citep{stein2025causalrivers}. Study 1 (P1--P3): the
benchmark's connected-subgraph sampler, seeds 7000--7029, East Germany and
Bavaria, $d\in\{5,10\}$, 120 subgraphs, 791 induced edges; block means at
12\,h and 24\,h; nominal Fisher-$z$. Study 2 (P4--P6): seeds 8000--8029, 120
subgraphs, 1{,}376 non-relative pairs; proxy $Z$ is the mean of standardised
stations of the region outside the subgraph and outside 2 undirected river
hops (429--615 stations); primary test on partial correlations with a residual
Bartlett $\neff$ ($K=100$ steps at 12\,h, 50 at 24\,h); median $\neff$ 139/113
(raw) and 231/211 (residual). Before conditioning, 0.98 (nominal) and 0.87
(effective) of non-relative pairs tested dependent. Median proxy
contamination by river relatives was 2.6\% (maximum 72\%), a case of the
contaminated-proxy term in \cref{prop:f3}. Post hoc and without verdict:
conditioning on $k=1,3,5,10$ principal components of the leave-out stations
gave dependent fractions $0.745,0.605,0.576,0.491$ (12\,h), with accuracy
$0.580$ and Spearman $0.16$ at $k=10$.

\textbf{Causal Chambers} \citep{gamella2025chambers}. Candidate selection
(recorded before analysis) compared Causal Chambers, CausalRivers with
measured precipitation (rejected: the driver reaches each gauge through its
own catchment filter, so by \cref{prop:f3} a scalar proxy cannot remove it),
a semi-synthetic cross-sectional manufacturing benchmark (rejected: not real,
not time series) and fMRI with physiological regressors (rejected: no ground
truth). Wind tunnel: runs \texttt{actuators\_random\_walk\_2}--\texttt{16} (run
1 excluded before analysis: truncated); 16 random variables, 26 edges, 63
non-relative pairs per run; block means over 10 rows, first differences,
standardisation, $n=999$; effective-$n$ Fisher-$z$ with Bartlett $K=20$.
Light tunnel: \texttt{actuators\_white}, 20 variables, 39 edges, ten segments
of 2{,}000 steps, levels. Predicted values come from the package's published
ground-truth graph with constant parameters removed.
Secondary analyses (no verdicts): nominal $n$ gives slightly more false
positives ($0.048$), as \cref{prop:nominal} predicts; levels instead of
increments gave $0.054$, contradicting our own secondary prediction that
levels of random walks would produce many spurious dependences; without block
means ($D=1$) relative recovery drops to $0.18$ and Spearman to $0.10$. In the
light tunnel, polariser angles act on two sensors through Malus' law, a
nonlinearity with near-zero linear correlation; as stated in the plan, these
relative pairs were never detected, without affecting orientation.

%=============================================================================
\section{The Real-Data Estimand $\tilde s$}
\label{app:estimand}

The theorems of \cref{sec:exact} concern $\srho$ of an unrolled DAG. The
real-data studies estimate something else: lag-0 dependences between the
observed series of one long record. This section defines the population
quantity they target and proves what it does and does not say.

\begin{definition}[estimand]\label{def:estimand}
Let $\GS=(V,\Es)$ be a summary graph (no loop assumption) and $O\subseteq V$
the observed nodes. For $v\in O$ let $\tilde\rho_O(v)=|\rel_S(v)\cap O|$
(relatives computed in the full graph $\GS$, so unobserved nodes act as latent
common ancestors). Let $E_O=\{(i,j)\in\Es:i,j\in O,\ i\ne j\}$ be the induced
non-self-loop edges and, for $E_O\neq\emptyset$,
$\tilde s(\GS,O)=\frac1{|E_O|}\sum_{(i,j)\in E_O}\srho\big(\tilde\rho_O(i),\tilde\rho_O(j)\big)$,
with the edge score of \cref{sec:setting}. Edges of the latent projection
that are not edges of $\GS$ (paths through unobserved nodes) are not scored.
\end{definition}

\textbf{Why this is the target.} Under the SVAR model of \cref{app:agg},
\cref{prop:zero} gives, for almost every parameter, $\Gamma(0)_{uv}\ne0$ iff
$\anc_S(u)\cap\anc_S(v)\ne\emptyset$ iff $v\in\rel_S(u)$, for every
observation scheme. So a perfect lag-0 dependence oracle on the observed
series declares exactly $\rel_S(v)\cap O$, and the relatives sort of the
observed counts, averaged over $E_O$, is $\tilde s$. (Lag-0 faithfulness is the
generic case of \cref{prop:zero}; it can fail for specific parameters,
e.g.\ through Malus-law nonlinearities.) Relation to \cref{thm:stat}: with
$O=V$, every window node of the stationary unrolled DAG has
$\rho_W(v(t))=T|\rel_S(v)|$, so $\tilde s(\GS,V)$ is the stationary value
computed over non-self-loop edges only.

\begin{proposition}[sign relation]\label{prop:estimand}
Let $E_O\ne\emptyset$ and $E_O^\uparrow=\{(i,j)\in E_O:\rel_S(i)\cap O\ne\rel_S(j)\cap O\}$.
(a) Every edge of $E_O$ scores $1/2$ or $1$, and
$\tilde s=\tfrac12+|E_O^\uparrow|/(2|E_O|)\ge\tfrac12$.
(b) $E_O^\uparrow\subseteq\Eu$. Hence $\tilde s>\tfrac12$ iff
$E_O^\uparrow\neq\emptyset$, and then some induced edge is an up-edge of $\GS$
($\Src(i)\ne\Src(j)$); if no induced edge is an up-edge, $\tilde s=\tfrac12$.
(c) $(i,j)\in E_O$ belongs to $E_O^\uparrow$ iff some observed $w$ lies in
$\rel_S(j)\setminus\rel_S(i)$, i.e.\ shares an ancestor with $j$ but none with
$i$; in particular an observed other parent $k$ of $j$ with
$\anc_S(k)\cap\anc_S(i)=\emptyset$ suffices.
(d) With $O=V$, $E_V^\uparrow=\Eu\setminus\{\text{self-loops}\}$, so
$\tilde s(\GS,V)-\tfrac12$ and $\srho^W-\tfrac12$ have the same sign.
(e) The converse of (b) fails: for the collider $a\to c\leftarrow b$ with loops
and $O=\{a,c\}$, the only induced edge $(a,c)$ is an up-edge, but
$\rel_S(a)\cap O=\rel_S(c)\cap O=\{a,c\}$, so $\tilde s=1/2$.
\lean{(a)--(d): \lname{relObs_subset_of_edge}, \lname{score_obsEdge}, \lname{sTilde_eq}, \lname{sTilde_ge_half}, \lname{upObs_subset_upEdges}, \lname{sTilde_gt_half_iff}, \lname{exists_upEdge_of_sTilde_gt_half}, \lname{sTilde_eq_half_of_no_up}, \lname{mem_upObs_of_witness}, \lname{upObs_univ} (file \lname{Observed.lean}); (e) by hand.}
\end{proposition}
\begin{proof}
(a) $i\in\anc_S(j)$ gives $\rel_S(i)\subseteq\rel_S(j)$
\citep[Thm.~1]{reisach2026topological}, hence the same after intersecting with
$O$; nested sets have equal or strictly ordered sizes. (b) If the observed parts
differ, the full sets differ, and by \cref{thm:stat} that is $\Src(i)\ne\Src(j)$.
(c) By (a) the observed sets are nested, so they differ iff the larger one has
an observed element outside the smaller; $w\in\rel_S(j)\setminus\rel_S(i)$ means
$\anc_S(w)\cap\anc_S(j)\ne\emptyset=\anc_S(w)\cap\anc_S(i)$; $k$ itself is such a
$w$. (d) is (b) with $O=V$ and $\Src(i)\ne\Src(j)$ iff $\rel_S(i)\ne\rel_S(j)$
(\cref{thm:stat}); $\srho^W>\tfrac12$ iff $\Eu\ne\emptyset$, and self-loops are
never in $\Eu$. (e) $\rel_S(a)=\{a,c\}$, $\rel_S(c)=\{a,b,c\}$; $b$ is unobserved.
\end{proof}

\textbf{What is and is not implied.} For river subgraphs (in-forests) an
induced edge $(i,j)$ counts iff an observed station lies upstream of $j$ on a
branch other than $i$'s (in an in-forest $\rel_S(v)=\anc_S(v)\cup\desc_S(v)$ and
$\desc_S(j)\subseteq\desc_S(i)$), so a confluence in the induced subgraph whose
other branch is observed is sufficient. $\tilde s$ is a statement about
population lag-0 relatives of a stationary process; it is not
$\srho$ of the finite unrolled DAG, of the induced subgraph's unrolled DAG, or
of the latent projection, and it says nothing about lagged dependence. A
deviation of the estimate from $\tilde s$ is a failure of estimation (F1--F3)
or of lag-0 faithfulness, not evidence against \cref{thm:edge,thm:limit}.

%=============================================================================
\section{Time-Stamped Analysis Plans: Process, Prior Exposure, Timeline, Deviations}
\label{app:prereg}

\textbf{Process.} The four analysis plans, all analyses and the results were
produced by AI coding agents (\cref{sec:real}, AI-use statement) within about
six hours of one day. Each plan was committed to version control before its
study's data were loaded for any dependence statistic in this project. Commit
timestamps are set by the committing client; no third-party registry (e.g.\
OSF) was used, and the hosting service's push records, the stronger evidence,
can be supplied to the area chair on request. The gap between plan and results
was short: about 3 min 41 s for plan 3 and 3 min 36 s for plan 4 (each
including a run of about two minutes). We therefore call these
\emph{internally time-stamped analysis plans}; earlier drafts called them
pre-registrations. Commit hashes are withheld for anonymity and will be
provided with the de-anonymised version.

\textbf{Prior exposure to CausalRivers (disclosure).} One day before plan 1, a
separate exploratory project in the same repository (a test of
contemporaneous-edge discovery on aggregated CausalRivers; commits dated
2026-11-03) had analysed CausalRivers with the same regions (East Germany,
Bavaria), at resolutions 6\,h, 12\,h and 24\,h (the studies here use 12\,h and
24\,h), and with the same data loaders and subgraph sampler that the plan-1 code
reuses (with new seeds). It ran correlation, VAR, PCMCI+, VAR-LiNGAM and
DYNOTEARS baselines and recorded a hidden-driver diagnostic: at 24\,h the VAR(1)
residual correlation among sibling stations was $0.51$, the same as among
unrelated non-adjacent stations in the same subgraph ($0.51$), and below
adjacent pairs ($0.65$); it concluded that shared regional forcing is large and
contemporaneous dependence is not specific to the river edges. This was known
to the agents (and recorded in the repository) before plan 1 was written. Plan
1 nevertheless predicted that the relatives sort would orient river edges
above chance; it failed. Because the strong shared dependence had been
observed beforehand, the F2 saturation account of that failure (and the F3
account of study 2) is \emph{explanatory}, not a prediction made blind to the
data. The Causal Chambers data and the E1/E2 held-out sets had not been
analysed before their plans.

Times below are commit timestamps relative to the first plan commit ($T_0$).

\begin{center}\small
\begin{tabular}{p{0.17\textwidth}p{0.70\textwidth}}
\toprule
time & event \\
\midrule
$T_0+0{:}00$ & plan 1 committed: targets T1--T3, predictions P1--P3, kill rule \\
$T_0+0{:}14$ & theory of \cref{sec:exact,sec:agg}, enumeration checks, operationalisation of P1--P3 (before data) \\
$T_0+0{:}17$ & study 1 results (P2, P3 fail; kill rule triggered) \\
$T_0+1{:}05$ & plan 2 committed: F1--F3 targets, Lean target, predictions P4--P6, kill rule \\
$T_0+1{:}18$--$1{:}31$ & Lean formalisation of \cref{sec:exact} and algebraic parts of \cref{sec:agg} \\
$T_0+1{:}38$ & finite-sample theory, simulations, operationalisation of P4--P6 (before data) \\
$T_0+1{:}44$ & study 2 results (P4, P6 fail; kill rule triggered) \\
$T_0+4{:}58$ & plan 3 committed: candidate scoping, graph-only predictions, P7--P9 (before any dependence statistic) \\
$T_0+5{:}02$ & study 3 results (P7, P8 pass; P9 fails; replication passes) \\
$T_0+5{:}14$--$5{:}24$ & Lean formalisation of the finite-sample lemmas; errata \\
$T_0+5{:}45$ & manuscript draft committed (with placeholders for E1/E2) \\
$T_0+5{:}50$ & plan 4 committed: E1 (test of the theorem against eight naive predictors), E2 (lagged relatives), theory addendum A1--A3 and B1--B3, graph-only predictions, design check; development runs on the study-3 data only; no dependence statistic of any held-out dataset computed \\
$T_0+5{:}54$ & E1/E2 results committed as planned (E1 passes 8/8 on the primary graph; Q7, Q8, Q10 pass, Q9 fails); out-of-sample F checks; run logs \\
\bottomrule
\end{tabular}
\end{center}

\textbf{Deviations (all recorded at the time).} (1) Plan 1 said
``theorem-predicted $\srho$'' without naming the version; the stationary
version was fixed in writing before data were loaded (the finite version would
predict $>1/2$ for every subgraph, including chains). (2) Self-loops are
excluded from the predicted average because the estimator cannot orient
them; this does not change which side of $1/2$ the prediction is on.
(3) Only 6 of 120 study-1 subgraphs lacked a confluence; the plan set no
minimum size for that group, and its trivial pass is not interpreted. (4) Plan
2 did not say which test defines ``dependent'' in P4; the effective-$n$ test
was fixed as primary before data, and the nominal test fails P4 as well.
(5) In study 3, a planned triad-reliability check was not applicable (the
measured driver reaches only four mutually related barometers) and is not
interpreted. (6) The P9 failure is close to the threshold and its CI contains
$0.3$; as fixed in the plan, the point estimate decides.
(7) Plan 4 (E1/E2): no threshold, unit, seed, $K$, $L$, $D$, graph or node set
changed after the results were seen. A stray cache file committed with plan 4 was removed in the results
commit (no content change). The F-check tolerance of $\pm0.03$ is the one fixed in plan 4; the citris
``predicted accuracy'' is undefined (no full-set edge of $G_{\mathrm{lin}}$ is predicted $1$). The
constant-$1/2$ comparison and the weak-component-size comparison are post hoc. The choice of $G_{\mathrm{lin}}$ as the primary graph was
fixed in plan 4, together with the secondary analysis on the published graph $G_D$, which is the one on which
the theorem fails against the sensor indicator on strictly ordered edges.



%=============================================================================
\section{Experiment E1: Discriminating Test Against Naive Predictors}
\label{app:e1}

\textbf{Status and provenance.} Specified in plan 4 (\cref{app:prereg}); the plan was
committed before any dependence statistic of any held-out dataset was computed. Done
before it: file hashes, column and row counts, regeneration of the
\texttt{citris\_1} input protocol from its published generator (it matches the recorded inputs exactly),
graph-only predictions, and development runs on the study-3 data only. All numbers below come from the
committed result files of the experiment folder (\texttt{results4.json}, \texttt{summary4.txt}); none is estimated or rounded beyond the printed digits.

\textbf{Why \texttt{citris\_\allowbreak 1}.} Every published Causal Chambers graph is bipartite except the pressure-control
configuration, whose only experiment varies no independent input, so we rejected it. The light-tunnel
\texttt{citris\_1} series (100{,}000 steps about 1.8\,s apart) is generated by a time-series structural model among the actuators:
red is a reflected random walk, green moves halfway towards red and blue a quarter of the way towards green (each plus a uniform step),
pol$_1$ is a random walk and pol$_2$ moves towards or away from pol$_1$ according to $\mathrm{sign}(\mathrm{blue}-\mathrm{red})$.
With the sensors this is a chain of length 3, red $\to$ green $\to$ blue $\to$ sensors. The secondary data are
snapshot (i.i.d.) structural models among the actuators
(\texttt{buchholz\_\allowbreak 1\_\allowbreak obs}; \texttt{scm\_\allowbreak \{1,2,4,5\}\_\allowbreak reference}, including a chain and a collider)
and the wind-tunnel regime-jump runs (bipartite; disagreements arise only from latent-node ties).

\textbf{Linear-zero lemma (the refinement).} Write $\mathcal C$ for the colour block (red, green, blue, ir$_1$,
vis$_1$, ir$_2$, vis$_2$, ir$_3$, vis$_3$, current) and $\mathcal Q$ for the polariser block (pol$_1$, pol$_2$, angle$_1$, angle$_2$).
\emph{B1 (proved for the actuators).} Conditionally on the flag and the colour paths, the law of the polariser
paths is invariant under negation, because every update (uniform draw on a symmetric interval, move towards or
away with a sign fixed by the colours, reflection at $\pm90$, rounding to $0.1$) commutes with negation.
\emph{B2 (assumed).} The first-two-layer sensors depend only on the colours and independent noise; ir$_3$ and vis$_3$
depend on the polarisers only through an even function of $\theta_1-\theta_2$ (Malus' law with aligned zero marks); each angle sensor is affine in its polariser plus
independent noise. This is a physical assumption that the design documents do not let us check.
\emph{B3.} Under B1 and B2 every variable of $\mathcal C$ has zero cross-covariance, at every pair of times, with every variable of
$\mathcal Q$ (a product of an even and an odd function of the polariser path has conditional mean zero). The six real edges
red $\to$ pol$_2$, blue $\to$ pol$_2$ and pol$_k\to$ ir$_3$, vis$_3$ therefore carry no linear dependence. The relatives that a
correlation-based estimator targets are those of the \emph{linear-dependence graph} $G_{\mathrm{lin}}$ (the design graph $G_D$ without these six edges),
which has two weak components, each with a single source (red; pol$_1$). By \cref{cor:statiff}, every edge inside a block has predicted value $1/2$ for every node set; the six cross-block
edges are compared by block size: red $\to$ pol$_2$ and blue $\to$ pol$_2$ are predicted $0$ (anti-sorted) and pol$_k\to$ ir$_3$, vis$_3$ are predicted $1$ for the full set.
Faithfulness inside the blocks is assumed, as in \cref{prop:zero}. A zero-mark offset or a non-affine angle sensor would break B2 and move the prediction towards $G_D$.

\textbf{Metric and tests.} For each edge instance (the full node set plus 20 subgraphs per unit) the observed value is
$o\in\{0,\frac12,1\}$ from the estimated counts, and the loss is $|o-p|$ for the predicted value $p$. For each competitor take the instances where it disagrees with the theorem;
$\Delta=$ loss(theory) $-$ loss(competitor), pooled, with a unit-level exact sign-flip test (1{,}024 flips for 10 units) and Holm correction over eight competitors, and a unit-bootstrap CI.
E1a (prerequisites): non-relatives declared $<0.5$ and relatives detected $\ge0.8$. E1b: $\Delta\le-0.10$ and Holm $p\le0.05$ for all eight. E1c: the same restricted to instances whose own pair is declared dependent ($\Delta<0$).
Separation was to be claimed iff E1a, E1b and E1c all pass on \texttt{citris\_\allowbreak 1}. Kill rules: K1 (E1a fails: no claim), K2 (a competitor is significantly better: no claim against it), K3 (E1b passes and E1c fails: attributed to power-induced ties).
None was triggered on the primary analysis.

\begin{table}[h]
\caption{E1 primary, \texttt{citris\_\allowbreak 1}, graph $G_{\mathrm{lin}}$, lag 0, 1{,}662 instances. $\Delta$ [95\% CI]; all rows: 10/10 units, Holm $p=0.0078$. ``Exact'': share of instances whose predicted value equals the observed one. ``Spearman'': agreement of the mean predicted edge value with observed accuracy over 200 subgraphs.}
\label{tab:e1prim}
\centering\small
\begin{tabular}{lcccc}
\toprule
predictor & $\Delta$ [CI] & $\Delta$, detected edges & exact & Spearman \\
\midrule
theorem & --- & --- & 0.992 & 0.962 \\
sensor indicator & $-0.502$ [$-0.507,-0.496$] & $-0.500$ & 0.203 & 0.408 \\
ancestors, full graph & $-0.515$ [$-0.519,-0.510$] & $-0.500$ & 0.147 & 0.334 \\
ancestors, subgraph & $-0.509$ [$-0.513,-0.505$] & $-0.500$ & 0.120 & 0.121 \\
in-degree, full graph & $-0.515$ [$-0.520,-0.511$] & $-0.500$ & 0.180 & 0.349 \\
in-degree, subgraph & $-0.510$ [$-0.514,-0.505$] & $-0.500$ & 0.143 & 0.105 \\
depth, full graph & $-0.515$ [$-0.519,-0.510$] & $-0.500$ & 0.147 & 0.334 \\
depth, subgraph & $-0.510$ [$-0.514,-0.505$] & $-0.500$ & 0.114 & 0.113 \\
every edge 1 & $-0.533$ [$-0.536,-0.530$] & $-0.500$ & 0.115 & undefined \\
\bottomrule
\end{tabular}
\end{table}

Observed values given the theorem's prediction: predicted $1/2$ (1{,}365 instances) observed exactly $1/2$ in $99.2\%$; predicted $1$ (187 instances,
polariser $\to$ ir$_3$, vis$_3$) observed $1$ in $99.5\%$; predicted $0$ (110 instances, red and blue $\to$ pol$_2$) observed $0$ in $99.1\%$. The detected-edge subset has 1{,}338 instances, all observed exactly $1/2$.
Prerequisites: non-relatives declared $0.000$ [0, 0] (40 pairs per unit), relatives detected $0.969$ [0.945, 0.992] (51 pairs per unit).

\textbf{Strict subset and the power safeguard.} On the 297 instances where the theorem predicts $0$ or $1$ it beats all eight ($\Delta=-0.539$ for the sensor indicator, $-0.948$ for ancestors and depth on the full graph, $-0.991$ for ``every edge 1'';
Holm $p=0.0078$, 10/10 units). There, an estimator without power would favour the competitors (the ``all counts tied'' baseline is $+0.443$ to $0$), so this part cannot come from power-induced ties.

\textbf{Published graph $G_D$ (secondary).} Applying the theorem to the unmodified design graph, relatives detection is $0.659$ [0.643, 0.675] (75 relative pairs per unit) and non-relatives declared $0.000$.
Overall the theorem still beats all eight ($\Delta=-0.323$ for the sensor indicator; $-0.424$ for ancestors, depth and ``every edge 1''; $-0.420$ and $-0.417$ for in-degree; 10/10 units, Holm $p=0.0078$), and its exact-match rate is $0.714$.
But on the 442 edges where $G_D$ predicts strict order ($1$), the theorem matches $0.222$ and the sensor indicator $0.319$; the sensor indicator is significantly better
($\Delta=+0.453$ [0.410, 0.492], theorem better in 0/10 units, reverse Holm $p=0.002$). A constant-$1/2$ rule (post hoc) matches $0.815$ of all instances, against $0.714$ for the theorem on $G_D$ and $0.992$ on $G_{\mathrm{lin}}$;
mean $|o-1/2|=0.092$ against $0.166$ for the theorem on $G_D$. The Spearman statistic is $0.341$ (theorem on $G_D$) against $0.408$ (sensor indicator). This is why the claim is restricted to the linearly dependent relatives.

\textbf{Lagged estimator on \texttt{citris\_\allowbreak 1} (secondary).} With $L=5$ and $G_{\mathrm{lin}}$: detection $0.976$ [0.953, 0.992], non-relatives declared $0.000$, 8/8 passed ($\Delta=-0.505$ to $-0.535$), theorem exact match $0.998$, Spearman $0.996$.

\begin{table}[h]
\caption{E1 on i.i.d.\ structural-model sets (no verdict; secondary). Competitors' exact-match range over the eight predictors. Units: \texttt{scm\_\allowbreak 1} 5, the others 10 each (45 pooled).}
\label{tab:e1iid}
\centering\footnotesize
\setlength{\tabcolsep}{3pt}
\begin{tabular}{lccccc}
\toprule
set & instances & theorem exact & competitors exact & $\Delta$ range & result \\
\midrule
\texttt{buchholz\_1\_obs} & 1{,}712 & 0.791 & 0.328--0.375 & $-0.499$ to $-0.496$ & pass 8/8 \\
\texttt{scm\_1\_reference} & 867 & 1.000 & 0.314--0.392 & $-0.500$ & not met$^\dagger$ \\
\texttt{scm\_2\_reference} & 1{,}650 & 0.812 & 0.324--0.408 & $-0.443$ to $-0.426$ & pass 8/8 \\
\texttt{scm\_4\_reference} (chain) & 2{,}267 & 1.000 & 0.000--0.088 & $-0.500$ & pass 8/8 \\
\texttt{scm\_5\_reference} (collider) & 2{,}238 & 1.000 & 0.464--0.495 & $-0.500$ & pass 8/8 \\
pooled (45 units) & 8{,}734 & 0.923 & 0.283--0.322 & $-0.489$ to $-0.487$ & pass 8/8, Holm $p=8\times10^{-5}$ \\
\bottomrule
\end{tabular}
\\[2pt]{\footnotesize $^\dagger$Five units: the smallest attainable Holm-adjusted $p$ is $8/32=0.25$.}
\end{table}

On the snapshot sets, relatives detected: $0.933$ (\texttt{buchholz\_\allowbreak 1\_\allowbreak obs}), $1.000$ (\texttt{scm\_\allowbreak 1}), $0.965$ (\texttt{scm\_\allowbreak 2}), $1.000$ (\texttt{scm\_\allowbreak 4}, \texttt{scm\_\allowbreak 5}); non-relatives declared $0.000$ wherever non-relative pairs exist (\texttt{scm\_\allowbreak 4} has none).
The chain (\texttt{scm\_\allowbreak 4}) gives $1/2$ on all 2{,}267 instances (observed exactly $1/2$ every time); the collider gives $1$ on edges into the collider and its sensors and $1/2$ elsewhere (all 2{,}238 instances exact).
Pooled, the theorem's strict predictions number 2{,}768 (observed value $1$ on 2{,}403, $1/2$ on 349, $0$ on 16: exact $0.868$); on them only the sensor indicator disagrees often (390 instances, $\Delta=-0.5$, 35/35 units) and the other competitors coincide with the theorem. The pooled Spearman is $0.913$ for the theorem and $-0.138$ for the sensor indicator.
Post hoc, a constant-$1/2$ rule matches $0.688$ of the pooled instances against $0.923$.

\textbf{Wind-tunnel regimes (tertiary, 64 units, no verdict).} The same test passes against all eight in each configuration. $D=10$, $L=0$: non-relatives declared $0.014$, relatives detected $0.592$, $\Delta=-0.203$ [$-0.227,-0.179$] in 63/64 units, exact match $0.539$ against $0.432$; detected-edge subset $\Delta=-0.169$.
$D=1$, $L=30$: $0.003$, $0.398$, $\Delta=-0.225$ [$-0.250,-0.197$] in 62/64 units, exact $0.443$ against $0.353$; detected-edge subset $\Delta=-0.134$. $D=1$, $L=0$: $\Delta=-0.218$ (64/64 units), exact $0.411$ against $0.326$; $D=10$, $L=3$: $\Delta=-0.218$, exact $0.533$ against $0.411$.
Detection of $0.33$--$0.59$ means that part of the advantage is power-induced, as plan 4 stated in advance. Post hoc, the constant-$1/2$ rule has mean absolute error $0.311$ against $0.363$ for the theorem at $D=1$, $L=30$ (and $0.313$ against $0.294$ at $D=10$, $L=0$), and matches $0.378$ ($0.375$) of the instances against $0.443$ ($0.539$).

\textbf{A trivial rule that E1 did not test (post hoc).} Each weak component of
$G_{\mathrm{lin}}$ has a single source (red; pol$_1$). By the proof of
\cref{cor:statiff}, every node of such a component is reached from that source,
so its relatives are the whole component, and for every node set the
predicted counts are the numbers of observed nodes in the endpoints' weak
components. The theorem's prediction therefore coincides on all 1{,}662
instances with the rule ``$1/2$ iff both endpoints lie in the same weak
component of $G_{\mathrm{lin}}$, otherwise compare the observed sizes of the two
components''. That rule was not among the eight competitors, so E1 cannot
separate the theory from it; the only data in this appendix where the source
structure matters beyond component membership are the snapshot sets with
several sources per component, notably the collider \texttt{scm\_\allowbreak 5}.

\textbf{What E1 does and does not show.} It shows that, on one multi-layer light-tunnel series, the estimated relatives counts are tied along the whole chain red $\to$ green $\to$ blue $\to$ sensors and between the polariser block and the angle sensors, and strictly ordered (including against the causal direction) on red, blue $\to$ pol$_2$ --- and that the source-component rule of \cref{thm:stat} and \cref{cor:statiff} predicts this while eight naive rules do not. It does not show that the theorem predicts relatives under the published graph; the faithfulness of the graph to linear dependence fails here by design. Most of the separation is in the ties, which naive rules cannot express, and on $G_{\mathrm{lin}}$ the prediction equals the weak-component-size rule. The evidence is one device in a favourable regime, and the wind tunnel, where detection is limited by F1, gives weaker support.

%-----------------------------------------------------------------------------
\section{Experiment E2: Lagged Relatives}
\label{app:e2}

\textbf{Theory (plan 4, addendum A).} For $L\in\{0,1,\dots\}$ let the lag-$L$ relatives of $u$ be $u$ and every $v$ with $\gamma_{uv}(h)\neq0$ for some $|h|\le L$.
\emph{A1 (proved).} Under the independent-innovations assumption behind \cref{prop:zero}, $\rel_0\subseteq\rel_L\subseteq\rel_\infty\subseteq\rel_S$; for a VAR model all are equal for almost every parameter value;
and \cref{thm:stat}, \cref{cor:statiff} hold verbatim with ``relatives'' read as $\rel_L$ under lag-$L$ faithfulness. Lag-0 faithfulness implies lag-$L$ faithfulness; the converse fails
(a white-noise source $s$ with an edge $s\to j$ and no self-loop at $s$ has $\gamma_{sj}(0)=0\ne\gamma_{js}(1)$).
\emph{Estimator.} This is a Bonferroni max-over-lags variant of the classical
cross-correlation tests of independence between two series
(\citealp{haugh1976checking} cross-correlates prewhitened residuals and uses a
portmanteau statistic; we use effective-$n$ Fisher-$z$ statistics without
prewhitening). Declare $u,v$ relatives iff $\max_{|h|\le L}|\atanh \hat r_{uv}(h)|\sqrt{\hat n_{\mathrm{eff},uv}(h)-3}>z_c(L)$, with $\hat n_{\mathrm{eff},uv}(h)=(n-|h|)/\hat B_{uv}$ and
$z_c(L)=\Phi^{-1}(1-\alpha/(2m(2L+1)))$, a Bonferroni correction over the $m$ pairs and the $2L+1$ lags; $L=0$ is exactly the study-3 test.
\emph{A2 (proved for finite $n$).} For independent stationary Gaussian $a,b$ the null variance of the lag-$h$ cross-covariance has the same bound at every $h$, so one Bartlett factor serves all lags.
\emph{A3 (deterministic part proved, constants sketched).} The recovery bound of F1 holds with $\rho_{\min}$ replaced by $\rho^*_{\min,L}=\min_{\text{relative pairs}}\max_{|h|\le L}|\rho_{ab}(h)|\ge\rho_{\min}$ (the gain) and a factor $2L+1$ in the failure probability and the Bonferroni count (the cost):
for 105 pairs at $\alpha=0.01$, $z_c$ rises from $3.90$ ($L=0$) to $4.35$ ($L=3$) and $4.80$ ($L=30$). A lagged correlation must therefore exceed the lag-0 correlation by about $1.12$ ($L=3$) or $1.23$ ($L=30$) times to help.

\textbf{Design.} Held-out data: the wind-tunnel regime-jump runs (\texttt{regime\_\allowbreak jumps\_\allowbreak single}, \texttt{regime\_\allowbreak jumps\_\allowbreak multi}), $2\times32=64$ regimes of $10{,}000$ steps, in the same archive as the study-3 data but never analysed in study 3.
Load\_\allowbreak in, load\_\allowbreak out, pot$_1$, pot$_2$ follow independent mirrored random walks within a regime; the hatch and sensor parameters are constant within a regime. The first 100 rows of each regime are dropped. Graph: 20 edges without the hatch (19 predicted $1$, pot$_1\to$ signal$_1$ predicted $1/2$); 54 non-relative and 51 relative pairs; 1{,}280 subgraphs (10 of size 5 and 10 of size 8 per unit). Primary: $D=1$ (first differences), $L=30$ (about 9\,s) against $L=0$; secondary: block means with $D=10$ and $L=3$ against $L=0$.
Success required Q7--Q10 all to pass at the primary configuration. Kill rules: Q7 fails (saturation) or Q10 fails (lags do not help detection, and the lagged estimator is not recommended). Plan 4 states that E2 is not a re-analysis of P9 and that P9's verdict is final.

\begin{table}[h]
\caption{E2, 64 held-out regimes, 95\% unit-bootstrap CIs. Q8: accuracy on the 1{,}216 instances predicted $1$. Q10: relatives detected, lagged minus lag 0 (non-relatives declared difference in the last column).}
\label{tab:e2}
\centering\tiny
\setlength{\tabcolsep}{3pt}
\begin{tabular}{lccccc}
\toprule
configuration & Q7 & Q8 & Q9 & Q10 & non-rel.\ diff. \\
\midrule
\textbf{primary $D{=}1$, $L{=}30$} & 0.003 [0.001, 0.005] pass & 0.553 [0.527, 0.579] pass & 0.243 [0.193, 0.294] \textbf{fail} & $+0.069$ [0.060, 0.079] pass & $+0.001$ [$-0.001$, 0.002] \\
secondary $D{=}10$, $L{=}3$ & 0.011 [0.008, 0.014] & 0.711 [0.692, 0.729] & 0.202 [0.152, 0.253] fail & $-0.007$ [$-0.012,-0.003$] fail & $-0.002$ [$-0.004$, $-0.0003$] \\
lag 0, $D{=}10$ (replication of study 3) & 0.014 [0.010, 0.018] & 0.742 [0.722, 0.760] & 0.202 [0.145, 0.259] fail & --- & --- \\
lag 0, $D{=}1$ & 0.002 [0.001, 0.004] & 0.552 [0.532, 0.570] & 0.120 [0.066, 0.176] fail & --- & --- \\
\bottomrule
\end{tabular}
\end{table}

Detection of the 51 relative pairs: $D=1$: $0.328$ [0.312, 0.345] at lag 0 and $0.398$ [0.383, 0.412] with $L=30$; $D=10$: $0.592$ [0.584, 0.599] and $0.585$ [0.576, 0.592] with $L=3$ (lags reduce detection: the Bonferroni cost exceeds the gain).
On \texttt{citris\_\allowbreak 1} with $L=5$ against $L=0$: $0.976$ against $0.969$, difference $+0.008$ [0, 0.024], so Q10 fails there; detection is already nearly saturated at lag 0.
\textbf{Verdict.} Q9 fails at the primary configuration, so E2 does not succeed; no kill rule is triggered (Q7 and Q10 pass). Lags raise detection at fine resolution without adding false positives, as A1 and A3 predict.
The study-3 estimator (lag 0) rerun on these held-out regimes fails the graded criterion again ($0.202$ at $D=10$, $0.120$ at $D=1$): the P9-type failure repeats on held-out data. P9's original verdict ($0.267$, failed) stands; E2 is not presented as rescuing it.
The development check of plan 4 showed that, on the study-3 data with the same estimator, drawing different subgraph seeds moves the P9 statistic from $0.267$ to $0.338$. We report this as variability of the statistic (about $0.07$ from the subgraph draw alone); it is a development run, not a verdict, and it does not change the P9 outcome.

\begin{table}[h]
\caption{Detection rate per edge over 64 regimes: share of regimes in which the pair was declared dependent.}
\label{tab:e2edges}
\centering\small
\begin{tabular}{lccc}
\toprule
edge & $D{=}1$, $L{=}0$ & $D{=}1$, $L{=}30$ & $D{=}10$, $L{=}0$ \\
\midrule
load\_in $\to$ rpm\_in & 0.859 & 0.875 & 1.000 \\
load\_in $\to$ rpm\_out & 0.000 & 0.172 & 0.094 \\
load\_in $\to$ current\_in & 1.000 & 1.000 & 1.000 \\
load\_in $\to$ current\_out & 0.000 & 0.000 & 0.047 \\
load\_out $\to$ rpm\_in & 0.000 & 0.016 & 0.000 \\
load\_out $\to$ rpm\_out & 0.844 & 0.969 & 1.000 \\
load\_out $\to$ current\_out & 1.000 & 1.000 & 1.000 \\
load\_out $\to$ current\_in & 0.000 & 0.000 & 0.000 \\
load\_in $\to$ pressure\_intake & 0.000 & 0.094 & 0.953 \\
load\_out $\to$ pressure\_intake & 0.000 & 0.031 & 0.906 \\
load\_in $\to$ pressure\_upwind & 0.562 & 0.984 & 1.000 \\
load\_out $\to$ pressure\_upwind & 0.016 & 0.812 & 1.000 \\
load\_in $\to$ pressure\_downwind & 0.734 & 0.953 & 1.000 \\
load\_out $\to$ pressure\_downwind & 0.750 & 1.000 & 1.000 \\
pot$_1\to$ signal$_1$ & 0.359 & 0.156 & 0.969 \\
pot$_1\to$ signal$_2$ & 0.266 & 0.109 & 0.969 \\
pot$_2\to$ signal$_2$ & 0.000 & 0.000 & 0.266 \\
pot$_1$, load\_in, load\_out $\to$ mic & 0.000 & 0.000 & 0.016 (pot$_1$), 0.000 \\
\bottomrule
\end{tabular}
\end{table}

Lags recover the pressure edges at fine resolution (load\_\allowbreak in $\to$ upwind pressure $0.562\to0.984$, load\_\allowbreak out $\to$ upwind pressure $0.016\to0.812$) but lose the pot edges to the larger threshold ($0.359\to0.156$ and $0.266\to0.109$); every microphone edge, pot$_2\to$ signal$_2$ and the cross-current edges stay undetected.
Even at its best ($D=10$, lag 0) wind-tunnel detection ($0.592$) is far below that of the light tunnel, so lags are not the main cause of the F1 shortfall.

\begin{table}[h]
\caption{Out-of-sample F-theory checks: the first half of each unit predicts the second half. Predicted detection is a bracket [$\bar p_{\max}$, $\bar p_{\mathrm{ind}}$] when lags are used. Pass if the observed value lies within 0.03 of the bracket (criterion fixed in plan 4).}
\label{tab:fcheck}
\centering\small
\setlength{\tabcolsep}{4pt}
\begin{tabular}{lcccc}
\toprule
configuration & predicted detection & observed & non-rel.\ declared, pred.\ / obs. & accuracy, pred.\ / obs. \\
\midrule
wind $D{=}1$, $L{=}0$ & 0.250 & 0.242 & 0.007 / 0.003 & 0.573 / 0.590 \\
wind $D{=}1$, $L{=}30$ & [0.344, 0.426] & 0.332 & [0.021, 0.074] / 0.003 & 0.504 / 0.506 \\
wind $D{=}10$, $L{=}0$ & 0.548 & 0.533 & 0.014 / 0.008 & 0.649 / 0.632 \\
wind $D{=}10$, $L{=}3$ & [0.536, 0.554] & 0.524 & [0.016, 0.027] / 0.008 & 0.628 / 0.618 \\
\texttt{citris\_1}, $L{=}0$ & 0.935 & 0.945 & 0.007 / 0.000 & undefined \\
\texttt{citris\_1}, $L{=}5$ & [0.940, 0.988] & 0.945 & [0.010, 0.045] / 0.000 & undefined \\
\bottomrule
\end{tabular}
\end{table}

All six checks pass. With lags the observed detection sits at or slightly below $\bar p_{\max}$, consistent with strongly correlated lag statistics; the independence formula overstates both detection and false positives, as the theory said it would.
For \texttt{citris\_\allowbreak 1} the predicted accuracy is undefined because no full-set edge of $G_{\mathrm{lin}}$ is predicted $1$ inside a block.

%=============================================================================
\section{Summary Graphs without Self-Loops}
\label{app:noloop}

Here $\GS=(V,\Es)$ is an arbitrary directed graph with $|\Es|\ge1$; no loop
assumption is made. $B_n$ is the boolean $n$-th power of the adjacency matrix
($B_0=I$), $\Lambda(u,v)=\{n:B_n[u,v]=1\}$ the set of lengths of walks from $u$
to $v$, and $u$ is \emph{fresh} if it has no incoming edge in $\Es$ (a self-loop
is an incoming edge). Status: \cref{thm:noloop} and \cref{cor:noloop} are
machine-checked; the rest of this section is proved on paper and checked by
enumeration, as marked.

\begin{lemma}[ancestors and roots without loops]\label{lem:noloopanc}
$u(a)\in\anc(v(t))$ iff $a\le t$ and $t-a\in\Lambda(u,v)$. A node $u(a)$ is a
root of $\GT$ iff $a=1$ or $u$ is fresh.
\lean{\lname{unrolled_reach_iff} (already free of the loop hypothesis).}
\end{lemma}
\begin{proof}
Every edge of $\GT$ raises time by one, so paths from $u(a)$ to $v(t)$ are the
lifts of walks of length exactly $t-a$, and every walk lifts. Without loops one
cannot pad. A node $u(a)$ with $a\ge2$ has a parent iff some edge $(w,u)$ exists.
\end{proof}

\begin{theorem}[edge value without loop assumption]\label{thm:noloop}
Let $(i,j)\in\Es$ and $e_t=(i(t),j(t{+}1))$, $1\le t<T$. Then
$\srho(e_t)\in\{\tfrac12,1\}$ independently of $T$, and $\srho(e_t)=1$ iff
$\mathrm{Strict}_{ij}(t)$ holds, where $\mathrm{Strict}_{ij}(t)$ is
\[
\exists u:\ t\in\Lambda(u,j),\ t-1\notin\Lambda(u,i)
\quad\text{or}\quad
\exists\text{ fresh }u,\ \exists n\in\{1,\dots,t\}:\ n\in\Lambda(u,j),\ n-1\notin\Lambda(u,i).
\]
\lean{\lname{isRel_unrolled_iff_NL}, \lname{reduce_ancestor}, \lname{relSet_unrolled_eq_iff_NL}, \lname{edge_score_NL}.}
\end{theorem}
\begin{proof}
As in \cref{thm:edge}, $\rel(i(t))\subseteq\rel(j(t{+}1))$, and the score is $1$
iff $\mathrm{roots}(j(t{+}1))\setminus\mathrm{roots}(i(t))\ne\emptyset$
\citep[Thm.~1, Lemma~1]{reisach2026topological}. By \cref{lem:noloopanc} an
element of the difference is $u(a)$ with ($a=1$ or $u$ fresh),
$n=t+1-a\in\Lambda(u,j)$ and $u(a)\notin\mathrm{roots}(i(t))$. The case $a=t+1$
would make $j$ fresh, impossible since $j$ has the parent $i$; so $a\le t$ and,
by the lemma again, $u(a)$ is not a root of $i(t)$'s ancestry iff
$n-1=t-a\notin\Lambda(u,i)$. The branch $a=1$ is $n=t$ with no condition on $u$;
the other is a fresh $u$ with $1\le n\le t$. (The Lean proof avoids Lemma~1: the
extra relative is $u(0)$ in the first case and a node of a fresh $u$ in the
second, and \lname{reduce_ancestor} pushes any common ancestor back along parents
to slice $0$ or to a fresh node.)
\end{proof}

\begin{corollary}\label{cor:noloop}
For every summary graph and $T\ge2$,
$\srho(\GT)=\tfrac12+\frac{1}{2(T-1)|\Es|}\sum_{(i,j)\in\Es}N_T(i,j)$ with
$N_T(i,j)=\#\{t<T:\mathrm{Strict}_{ij}(t)\}$; in particular $\srho(\GT)\ge1/2$.
With a loop at every node no node is fresh, $\Lambda(u,j)=[d(u,j),\infty)$, and
$\mathrm{Strict}_{ij}(t)$ becomes $\exists u:d(u,j)\le t\le d(u,i)$, i.e.\
\cref{thm:edge}.
\lean{\lname{sRho_unrolled_eq_NL}, \lname{sRho_unrolled_ge_half_NL}.}
\end{corollary}
Exhaustive check of \cref{thm:noloop} against brute-force relatives: $p=2,3$
($T=2,\dots,9$), all $65{,}535$ graphs with $p=4$ ($T=3,5,8$), and random graphs
with $p=5,6$; $0$ mismatches in $120$, $3{,}577$, $196{,}605$, $8{,}928$ and
$4{,}494$ (graph,~$T$) cases (enumerated, not part of the Lean proof).

\begin{proposition}[limit exists and is rational; paper proof, enumerated]\label{prop:noloopperiod}
$(B_n)$ is eventually periodic: there are $n_0,\pi\ge1$ with $B_{n+\pi}=B_n$
for $n\ge n_0$. Hence $L=\lim_T\srho(\GT)=\tfrac12+\frac1{2|\Es|}\sum_e\delta_e$
with $\delta_e=1$ if some fresh $u$ has $n\ge1$, $n\in\Lambda(u,j)$,
$n-1\notin\Lambda(u,i)$, and otherwise $\delta_e$ the density, over one period
of $t>n_0$, of $t$ with some $u$, $B_t[u,j]=1$, $B_{t-1}[u,i]=0$; and
$|\srho(\GT)-L|\le(n_0+\pi+1)/(2(T-1))$.
\end{proposition}
\begin{proof}
$B_{n+1}$ is a function of $B_n$ and there are finitely many boolean matrices.
A fresh witness is monotone in $t$ and has its first $n$ at most $n_0+1+\pi$
(the set $\Lambda(u,j)\setminus(\Lambda(u,i)+1)$ is eventually periodic), so
such an edge is strict for all $t\ge n_0+1+\pi$. Without a witness,
$\mathrm{Strict}_{ij}(t)$ depends for $t>n_0$ only on $t\bmod\pi$. Counting
strict $t<T$ gives the error bound; average over edges.
\end{proof}
Checked on $1{,}115$ graphs (all $15+511$ graphs with $p=2,3$, $390$ random with
$p=4$, $199$ random with $p=5,6$): the exact value at $T=600$ is within the
bound and $(L-\tfrac12)\,2|\Es|\,\pi$ is an integer in every case.

\begin{proposition}[edges inside an SCC; paper proof, enumerated]\label{prop:noloopscc}
Let $(i,j)\in\Es$ with $j$ reaching $i$ and no fresh witness. Then
$\mathrm{Strict}_{ij}(t)$ fails for every $t\ge n_0+1$.
\end{proposition}
\begin{proof}
Fix $u$; put $f(n)=B_n[u,j]$, $h(n)=B_{n-1}[u,i]$. The edge gives $h(n)\le f(n)$.
If $m$ is the length of a walk $j\to i$ and $c=m+1$, then $f(n)\le h(n+c)$. Take
$k\ge1$ with $kc$ a multiple of $\pi$. For $n\ge n_0+1$,
$h(n)\le f(n)\le h(n+c)\le\dots\le h(n+kc)=h(n)$, so $f(n)=h(n)$ for every $u$.
\end{proof}

\begin{proposition}[crossing edge into a cyclic node]\label{prop:noloopcross}
If $(i,j)$ is crossing and $j$ lies on a cycle with period $g$ (gcd of the
closed-walk lengths at $j$), then $\mathrm{Strict}_{ij}(t)$ holds for every
$t\in\Lambda(j,j)$, which contains all large multiples of $g$, so
$\delta_e\ge1/g$, and $\delta_e=1$ if $g=1$. A fresh parent $u\ne i$ of $j$ makes
every layer strict.
\lean{\lname{strictNL_of_closed_walk}, \lname{strictNL_of_fresh_parent}, \lname{sRho_unrolled_eq_one_NL}.}
\end{proposition}
\begin{proof}
Take $u=j$: $t\in\Lambda(j,j)$ and $t-1\notin\Lambda(j,i)=\emptyset$. Closed-walk
lengths at $j$ form a numerical semigroup in $g\mathbb Z$, which contains all
large multiples of $g$.
\end{proof}

\begin{corollary}[loop-free versions of the loop-dependent statements]\label{cor:noloopmain}
(a) If every weakly connected component of $\GS$ is strongly connected (any
loop pattern, any periods), then $L=1/2$ and $\srho(\GT)-\tfrac12\le n_0/(2(T-1))$
(paper proof from \cref{prop:noloopscc}; Reisach et al.'s Theorem~4 without
loops). (b) The converse and the closed form are false without loops
(\cref{tab:noloop}). (c) If every crossing edge has a head on a cycle, then
$L>1/2$ iff $\Ex\ne\emptyset$.
\end{corollary}
\begin{table}[h]
\caption{Loop-free examples ($L$: finite limit, $S$: stationary). Lean: \lname{edgeNoLoops_sRho}, \lname{colliderNL_sRho}, \lname{edgeLoopHead_sRho}; the others are computed values (example scripts in the supplement).}
\label{tab:noloop}
\centering\small
\begin{tabular}{lcc}
\toprule
summary graph & without loops & with loops \\
\midrule
single edge $0\to1$ & $\srho(\GT)=1/2$ for all $T$ & $L=2/3$, $S=1/2$ \\
chain $0\to1\to2$ & $L=1/2$ & $L=0.7$ \\
2-cycle $\{0,1\}$ plus $0\to2$ & $L=1/2$ & loop at $2$ only: $L=5/8$ \\
collider & $L=S=1$ & $L=S=0.7$ \\
edge, loop at head only & $L=S=1$ & -- \\
\bottomrule
\end{tabular}
\end{table}
For an acyclic loop-free graph, $L(u,v)$ is
finite, $B_n=0$ for $n\ge p$, and $L=\tfrac12+\#\{e:\text{fresh witness}\}/(2|\Es|)$
(paper proof); for a loop-free in-forest an edge $(i,j)$ has a fresh witness iff
$j$ has another parent, so $L=\tfrac12+(m-n_1)/(2m)$ with $m=|\Es|$ and $n_1$ the
number of nodes with one parent, which equals the stationary value below. On
$362$ random acyclic loop-free graphs $L=S$ in all cases, and the in-forest
formula holds on $298$ in-forests.

\begin{theorem}[stationary value without loops; (a), (c) proved, (b) sketched]\label{thm:noloopstat}
Let $\GS$ be arbitrary, $G_{\mathbb Z}$ the infinite-past unrolled DAG on
$V\times\mathbb Z$, and $\Delta(v,w)=\{m-n:\exists u,\ n\in\Lambda(u,v),\ m\in\Lambda(u,w)\}$.
Then $\rel_{G_\mathbb Z}(v(t))=\{w(t+d):d\in\Delta(v,w)\}$ (proved), the edge
$(i(t),j(t{+}1))$ is strict iff $\Delta(j,w)+1\ne\Delta(i,w)$ for some $w$
(proved; always $\Delta(i,w)\subseteq\Delta(j,w)+1$), and
$S^{\mathbb Z}=\tfrac12+|E^{\mathbb Z}_\uparrow|/(2|\Es|)$. (a) Outside
$|d|<n_0+\pi$, $\Delta(v,w)$ is $\pi$-periodic, so $E^{\mathbb Z}_\uparrow$ is detected
with witnesses of $|d|\le n_0+2\pi$ (proved). (b) In a finite window after a
burn-in $B\ge 2n_0+3\pi+1$ the windowed value equals $S^{\mathbb Z}+O(K_0/T)$,
$K_0=n_0+2\pi$ (\emph{sketched}). (c) $S^{\mathbb Z}=1/2$ whenever every weak
component is strongly connected, whatever its periods (a node colouring
$\mathrm{lab}$ with $\Delta(v,w)=\mathrm{lab}(w)-\mathrm{lab}(v)+g\mathbb Z$)
(proved). With loops everywhere $\Delta$ is $\mathbb Z$ or empty and this is
\cref{thm:stat}.
\end{theorem}
Brute-force relatives in $G_{B+T}$ ($B=40$, $T=90$, edges at distance $\ge30$
from the window ends) agree with $E_\uparrow^{\mathbb Z}$ on every edge of all
$15+511$ graphs with $p=2,3$ and $248$ random graphs with $p=4$ (enumerated).
Not formalised in Lean: periodicity, \cref{prop:noloopscc}, \cref{cor:noloopmain}(a),
\cref{thm:noloopstat}.

%=============================================================================
\section{Lagged Unrollings}
\label{app:lags}

\Cref{thm:edge,thm:limit,thm:stat} and their Lean proofs are for VAR(1)-shaped
unrollings: every summary edge acts at lag one. Here every non-loop summary
edge $e=(i,j)$ carries a lag $\tau_e\in\{1,\dots,\tau_{\max}\}$ (a pair may carry
several lags, each a separate edge), self-loops act at lag one, and $\GT$ has
the edges $(i(t),j(t{+}\tau_e))$, $1\le t\le T-\tau_e$. Let $d_\tau(u,v)$ be the
smallest total lag of a path from $u$ to $v$ ($0$ if $u=v$, $\infty$ if none);
$d_\tau(u,v)\le(p-1)\tau_{\max}$ when finite. Status: paper proof, not
machine-checked and not enumerated. Contemporaneous (lag-0) edges are excluded:
they break the time-monotonicity used below.

\begin{proposition}[lag-weighted version]\label{prop:lags}
(a) $u(s)\in\anc(v(t))$ iff $d_\tau(u,v)\le t-s$. (b) The copy
$(i(t),j(t{+}\tau_e))$ scores $1$ iff some $u$ has
$d_\tau(u,j)-\tau_e+1\le t\le d_\tau(u,i)$, and $1/2$ otherwise. (c) Crossing
edges score $1$ in every copy; an edge inside an SCC has at most
$(p-1)\tau_{\max}$ strict copies, so
$\lim_{T\to\infty}\srho(\GT)=\tfrac12+|\Ex|/(2|\Es|)$ with $|\srho(\GT)-L|=O(p\tau_{\max}/T)$.
(d) With a latent burn-in $B\ge(p-1)\tau_{\max}$ (or an infinite past),
\cref{thm:stat} holds verbatim.
\end{proposition}
\begin{proof}
(a) Every edge raises time by its lag, so a path from $u(s)$ to $v(t)$ projects
to a walk of total lag $t-s$; conversely a path of total lag $d_\tau(u,v)$ can be
padded by lag-one self-loops at $u$. (b) As in \cref{lem:roots}, the roots of
$v(t)$ are $\{u(1):d_\tau(u,v)\le t-1\}$ (every $u(s)$, $s\ge2$, has the parent
$u(s-1)$), and $i(t)$ is a parent of $j(t{+}\tau_e)$, so the score is $1$ iff
$j(t{+}\tau_e)$ has a root $u(1)$ with $d_\tau(u,j)\le t+\tau_e-1$ that $i(t)$
lacks, $d_\tau(u,i)>t-1$. (c) For a crossing edge take $u=j$:
$d_\tau(j,j)=0\le t+\tau_e-1$ and $d_\tau(j,i)=\infty$. Inside an SCC,
$d_\tau(u,i)<\infty$ iff $d_\tau(u,j)<\infty$, so a strict copy needs
$t\le d_\tau(u,i)\le(p-1)\tau_{\max}$. Each edge has $T-\tau_e$ copies; average.
(d) By (a), $w(s)\in\rel(v(t))$ iff some $u$ has $d_\tau(u,v)\le t-1$ and
$d_\tau(u,w)\le s-1$; for window nodes $t-1,s-1\ge B\ge(p-1)\tau_{\max}$ this is
$\anc_S(v)\cap\anc_S(w)\ne\emptyset$, and the proof of \cref{thm:stat} applies.
\end{proof}

%=============================================================================
\section{Random Summary Graphs: Derivations and Checks}
\label{app:er}

All graphs here are a non-loop random graph plus a self-loop at every node, so
$|\Es|=p+m$ with $m$ the number of non-loop edges, and for acyclic graphs
$|\Ex|=m$.

\textbf{In-forests (proved identities).} Let the non-loop part be an in-forest
with $m$ edges and $n_1$ nodes with exactly one parent. By
\cref{cor:iff,cor:statiff},
$L=\tfrac12+\frac{m}{2(m+p)}$, $S=\tfrac12+\frac{m-n_1}{2(m+p)}$,
$L-S=\frac{n_1}{2(m+p)}$ (the last checked to machine precision on $18$ random
forests). For a sequence of in-forests with parent-count distribution $\pi$ and
mean $\mu$: $L\to\tfrac12+\frac\mu{2(1+\mu)}$ and
$S\to\tfrac12+\frac{\mu-\pi_1}{2(1+\mu)}$. Exact computation at $p=20{,}000$
(6 seeds each): uniform recursive in-tree ($\pi_1=1/4$, $S\to11/16=0.6875$,
measured $0.68739\pm0.00026$); preferential attachment (measured $\pi_1=0.1656$,
$S=0.70859\pm0.00028$); uniform random labelled tree rooted and oriented to the
root (Poisson(1) parents, $\pi_1\to e^{-1}$, $S\to\tfrac12+(1-e^{-1})/4=0.65803$,
measured $0.65783\pm0.00020$). In all three $L=0.74999$.

\textbf{Directed ER: summary.} $G(p,c/p)$, $m/p\to c$. The finite limit needs
the fraction of non-loop edges inside an SCC ($\to x^2$, the two endpoints lie in
the giant SCC iff the ancestor tree of $i$ and the descendant tree of $j$ are
infinite); the stationary value needs the fraction of up-edges
($\to F_\uparrow$). Both are proved in \cref{app:erproof} below. Large $c$:
$1-x\sim e^{-c}$, $L-\tfrac12\sim e^{-c}$, $S-\tfrac12\sim e^{-c}/2$.

\textbf{Layered DAG (sketch; $L$ proved given edge concentration).} $M$ layers of
width $w$, each node of layer $l+1$ receiving an edge from each node of layer $l$
with probability $c/w$; as $w\to\infty$ with $\mu=c(M-1)/M$,
$L\to\tfrac12+\frac\mu{2(1+\mu)}$ and
$S\to\tfrac12+\frac{\mu(1-e^{-c})}{2(1+\mu)}$. For fixed $w$ and deep $M$ the
stationary value falls below this (observed $0.7824$ against $0.7878$ at
$w=50$, $M=400$, $c=2$); this correction is not derived.

\begin{table}[h]
\caption{Exact computation (means $\pm$ s.d.\ over seeds; SCCs by scipy and source sets by bitsets) against the limits of \cref{prop:er}, proved in \cref{app:erproof}.}
\label{tab:er}
\centering\small
\begin{tabular}{rrcccc}
\toprule
$c$ & $p$ & $L$ exact & $L$ formula & $S$ exact & $S$ formula \\
\midrule
0.5 & 20000 & $0.6666\pm0.0011$ & $0.6667$ & $0.5658\pm0.0014$ & $0.5656$ \\
1.0 & 20000 & $0.7486\pm0.0008$ & $0.7500$ & $0.6557\pm0.0013$ & $0.6580$ \\
1.5 & 50000 & $0.6975\pm0.0020$ & $0.6981$ & $0.6159\pm0.0018$ & $0.6163$ \\
2.0 & 50000 & $0.6214\pm0.0022$ & $0.6217$ & $0.5629\pm0.0014$ & $0.5632$ \\
3.0 & 100000 & $0.5434\pm0.0010$ & $0.5433$ & $0.5215\pm0.0006$ & $0.5214$ \\
5.0 & 20000 & $0.5059\pm0.0005$ & $0.5058$ & $0.5028\pm0.0003$ & $0.5029$ \\
\bottomrule
\end{tabular}
\end{table}
$L$ at $c=1$ converges as $p$ grows: $0.7466$ ($p=2000$), $0.7486$ ($2\cdot10^4$),
$0.74976$ ($2\cdot10^5$), $0.74992$ ($10^6$); at $c=2$, $0.62162\pm0.00023$ at
$p=10^6$ against $0.62170$. The giant-SCC fraction matches $x^2$ ($0.6353$ against
$0.6349$ at $c=2$). The second fixed point matters: with $d=1$ the formula would
give $S=0.6312$, $0.5766$, $0.5246$ at $c=1.5,2,3$, outside the error bars,
whereas the fixed points $d=0.485$, $0.221$, $0.061$ give $0.6163$, $0.5632$,
$0.5214$. Further exact rows: $c=1.5$, $p=10^5$, 32 graphs, $L=0.69819\pm0.00024$
(formula $0.69810$); $c=2$, $p=3.2\cdot10^5$, $L=0.62170\pm0.00012$ ($0.62170$).
The Barak--Erd\H{o}s phase-transition paper cited by \citet{reisach2026topological}
concerns monotone-path reachability at the dense scale $\log n/n$ and is not
used here.

\subsection{Directed Erd\H{o}s--R\'enyi summary graphs: full proof of \cref{prop:er}}
\label{app:erproof}

\textbf{Status.} Everything in this subsection is proved on paper. The only
inputs used without proof are textbook facts (F1)--(F3) below, checked against
the author's online version of \citet{hofstad2016random} (theorem numbers are
those of that version and may differ in print). Lemma~\ref{lem:er-local} is the
standard local-coupling argument for sparse random graphs and is written
concisely. The deterministic witness lemmas are machine-checked; no
probabilistic step is. No human probabilist has yet checked this subsection
(see the AI-use statement).

\textbf{Limits and their order.} For a fixed summary graph, \cref{thm:limit}
gives $\lim_{T\to\infty}\srho(\GT)=L(\GS)=\tfrac12+|\Ex|/(2|\Es|)$, and
\cref{thm:stat} gives $\srho^W=\tfrac12+|\Eu|/(2|\Es|)$ for every $T$. For
$\GS\sim D(p,c/p)$ we take the limit $T\to\infty$ at fixed $p$ first (this
defines $L(\GS)$) and then let $p\to\infty$; for the stationary value no
$T$-limit is needed. The regime of fixed $T$ with $p\to\infty$ is not covered
(the bound of \cref{thm:limit} is vacuous for $T\le p$).

\textbf{Model and notation.} $D(p,c/p)$: each ordered pair $(i,j)$, $i\ne j$,
is a non-loop edge independently with probability $q=c/p$, and every node has
a self-loop (loops change neither reachability nor $\Src$). $m$ is the number of
non-loop edges, $|\Es|=p+m$, $N$ the number of non-loop edges inside an SCC, so
$|\Ex|=m-N$. $\mathrm{out}(v)$, $\mathrm{in}(v)$ are the sets of nodes
reachable from $v$, resp.\ reaching $v$ ($v$ included), $T^{\mathrm{out}}_v$,
$T^{\mathrm{in}}_v$ their sizes; $v\leadsto w$ means $w\in\mathrm{out}(v)$.
$x=x(c)$ is the largest solution of $x=1-e^{-cx}$ (the survival probability of
a Poisson$(c)$ Galton--Watson tree), $\theta_K=\theta_K(c)$ the probability that
its total progeny is at least $K$ ($\theta_K\downarrow x$), and $d$ the smallest
fixed point in $[0,1]$ of $f(d)=1+e^{-c}-e^{-c(1-x)d}$. A node is
\emph{in-isolated} if it has no non-loop parent. $\mathcal G$ denotes the
\emph{edge-rooted Poisson tree}: an edge $(i,j)$; $i$ has Poisson$(c)$ parents
and Poisson$(c)$ further children, $j$ Poisson$(c)$ further parents and
Poisson$(c)$ children; every other node has Poisson$(c)$ parents and
Poisson$(c)$ children besides the edge through which it was reached; all
independent.

\textbf{Textbook facts.} (F1) Poisson coupling \citep[Thm.~2.10]{hofstad2016random}:
$d_{\mathrm{TV}}(\mathrm{Bin}(n,q),\mathrm{Poi}(c))\le nq^2+|nq-c|$.
(F2) Chernoff: for a sum $X$ of independent Bernoulli variables with mean
$\mu$ and $\delta\in(0,1)$, $P(X\le(1-\delta)\mu)\le e^{-\delta^2\mu/2}$. (From
$Ee^{-tX}\le\exp(\mu(e^{-t}-1))$ with $t=-\log(1-\delta)$, Markov's inequality,
and $(1-\delta)\log(1-\delta)\ge-\delta+\delta^2/2$; the upper-tail version
$P(X\ge\mu+s)\le e^{-s}$ for $s\ge e^2\mu$ follows in the same way.)
(F3) Galton--Watson facts \citep[Thms.~3.1, 3.16]{hofstad2016random}: the
extinction probability is the smallest root of $\eta=e^{c(\eta-1)}$, so the
survival probability is $x(c)$, positive iff $c>1$.

\textbf{Deterministic lemmas} (machine-checked, \lname{ErLemmas.lean}; no loop
hypothesis). A node $z$ is a \emph{source node} if every ancestor of $z$ is
reachable from $z$. By \cref{thm:stat}, $(i,j)\in\Eu$ iff some source node
reaches $j$ but not $i$.
\begin{itemize}
\item[(W1)] If $(i,j)\in\Eu$, there are $k\notin\{i,j\}$ with $k\to j$ and a
source node $z$ with $z\leadsto k$, $z\not\leadsto i$ (\lname{upEdge_witness}).
\item[(W2)] If $i\to j$, $k\to j$, and a source node reaches $k$ but not $i$,
then $(i,j)\in\Eu$ (\lname{mem_upEdges_of_witness}).
\item[(W3)] If $i\to j$, $k\to j$ and $\mathrm{in}(k)\cap\mathrm{in}(i)=\emptyset$,
then $(i,j)\in\Eu$ (\lname{mem_upEdges_of_disjoint_anc}).
\item[(W5)] A non-loop edge $(i,j)$ lies inside an SCC iff $j\leadsto i$. The
event $\{j\leadsto i\}$ and the set $\mathrm{out}(j)$ do not change when the
edge $(i,j)$ is added or removed (a simple path from $j$ to $i$ does not use
it). (Trivial; not formalised.)
\end{itemize}

\textbf{Explorations and forced edges.} Forward exploration of $\mathrm{out}(v)$
keeps explored, active and unseen sets $X_t,A_t,U_t$, $a_t=|A_t|$,
$\mathrm{disc}_t=t+a_t$, $X_0=\emptyset$, $A_0=\{v\}$. Step $t+1$ takes the
oldest active node $w$, reveals the pairs $(w,u)$, $u\in U_t$, and makes the
$\xi_{t+1}$ heads of present pairs active; it stops at $T=T^{\mathrm{out}}_v$.
Backward exploration of $\mathrm{in}(v)$ reveals the pairs $(u,w)$ instead.
Let $F$ be a set of at most three ordered pairs and $G^F$ the graph in which the
pairs of $F$ are present and all other pairs are independent Be$(q)$.
\emph{Bookkeeping.} In a forward exploration a pair is revealed only when its
tail is explored, and each node is explored once, so each pair is revealed at
most once; the pairs revealed at step $t+1$ are disjoint from those revealed
before and from all pairs whose tail is not in $X_{t+1}$. Hence, conditionally
on the history $\mathcal H_t$, the non-forced pairs revealed at step $t+1$ are
independent Be$(q)$, and
$\xi_{t+1}=\#\{u\in U_t:(w,u)\in F\}+\mathrm{Bin}(|U_t|-f_w,q)$ with
$f_w\le3$ the number of forced pairs from $w$ into $U_t$. Thus, given
$\mathcal H_t$, $\xi_{t+1}$ stochastically dominates $\mathrm{Bin}(|U_t|-3,q)$
and is dominated by $3+\mathrm{Bin}(p,q)$. The same holds for backward
explorations (with heads in place of tails), and, by the same argument, for a
backward exploration run after a forward one, conditionally on the forward
history, as long as it never discovers a node that the forward exploration
discovered (only pairs with tail outside $X_f$ and head explored backward are
revealed, and none of them was revealed forward).

Fix $c>1$ and $\varepsilon\in(0,1/5)$ with $\beta=c(1-4\varepsilon)>1$; put
$n'=\lceil(1-4\varepsilon)p\rceil$, $\kappa=(\beta-1)^2/(8\beta)$,
$\delta=(\beta-1)\varepsilon/(\beta+1)$, $t_0=C_3\log p$ with $C_3\kappa\ge8$.
Call an exploration \emph{admissible} if, conditionally on its history, every
offspring count $\xi_{t+1}$ taken while $\mathrm{disc}_t<\varepsilon p$
dominates $\mathrm{Bin}(n',q)$ and every offspring count is dominated by
$3+\mathrm{Bin}(p,q)$. By the bookkeeping, forward and backward explorations in
$G^F$ are admissible for large $p$ (since $|U_t|-3>(1-\varepsilon)p-3\ge n'$),
and so is the backward exploration of Lemma~\ref{lem:er-l2} below.

\begin{lemma}[dichotomy]\label{lem:er-l1}
There is $C$ (depending on $c,\varepsilon$) such that for large $p$ every
admissible exploration with final size $T$ satisfies
(a) $P(K\le T<\varepsilon p)\le Ce^{-\kappa K}$ for all $K\ge1$;
(b) $P\big(T\ge t_0$ and [$\mathrm{disc}_t<\varepsilon p$ for all $t$, or
$a_\tau<\delta p$ at $\tau=\min\{t:\mathrm{disc}_t\ge\varepsilon p\}$]$\big)\le Cp^{-8}$.
Consequently, with probability at least $1-Cp^{-7}$, every node $v$ of $G$ has, in
both directions, $T_v<t_0$ or ($T_v\ge\varepsilon p$ and $a_\tau\ge\delta p$).
\end{lemma}
\begin{proof}
Let $\xi'_t$ be independent $\mathrm{Bin}(n',q)$ variables coupled so that
$\xi_t\ge\xi'_t$ for $t\le\tau$ (admissibility; after $\tau$ or after the
exploration stops, $\xi'_t$ uses external randomness), and
$W_t=1+\sum_{s\le t}(\xi'_s-1)$. By induction $a_t\ge W_t$ for
$t\le\min(T,\tau)$. Since $tn'q\ge t\beta$ and
$(\beta+1)/2=(1-\frac{\beta-1}{2\beta})\beta$, (F2) gives
$P(W_t\le(\beta-1)t/2)\le P(\mathrm{Bin}(tn',q)\le t(\beta+1)/2)\le e^{-\kappa t}$,
so $\mathrm{Bad}(t_1)=\{\exists t\ge t_1:W_t\le(\beta-1)t/2\}$ has probability at
most $e^{-\kappa t_1}/(1-e^{-\kappa})$.
(a) If $T<\varepsilon p$ then $\tau=\infty$, so $0=a_T\ge W_T$, which is
$\mathrm{Bad}(K)$ when $T\ge K$.
(b) On $\mathrm{Bad}(t_0)^c$ with $T\ge t_0$: for $t_0\le t\le\min(T,\tau)$,
$a_t\ge W_t>0$, so the exploration cannot stop before $\tau$, hence
$\tau<\infty$. At $\tau$, $a_\tau+\tau\ge\varepsilon p$, and if $\tau\ge t_0$
also $a_\tau>(\beta-1)\tau/2$; minimising $\max((\beta-1)\tau/2,\varepsilon p-\tau)$
over $\tau$ gives $a_\tau\ge\delta p$; if $\tau<t_0$, $a_\tau\ge\varepsilon p-t_0\ge\delta p$.
$e^{-\kappa t_0}\le p^{-8}$. The consequence is a union bound over $p$ nodes and
two directions.
\end{proof}

\begin{lemma}[large out-sets reach large in-sets]\label{lem:er-l2}
With probability at least $1-Cp^{-5}$: for all $v,w$, $T^{\mathrm{out}}_v\ge t_0$
and $T^{\mathrm{in}}_w\ge t_0$ imply $v\leadsto w$.
\end{lemma}
\begin{proof}
Fix $(v,w)$. Explore $\mathrm{out}(v)$ forward up to $\tau_f$ (history $\mathcal H_f$,
discovered set $R_f=X_f\cup A_f$). Since every $\xi$ is dominated by
$3+\mathrm{Bin}(p,q)$, $P(\max_t\xi_t\ge\varepsilon p)\le e^{-\varepsilon p}$ for
large $p$ (F2, upper tail), so $|R_f|\le2\varepsilon p$ except on that event.
Then explore $\mathrm{in}(w)$ backward, stopping at $\tau_b$ or when a node of
$R_f\subseteq\mathrm{out}(v)$ is discovered; in the latter case $v\leadsto w$.
Otherwise the backward exploration is the true exploration of $\mathrm{in}(w)$,
its unseen set has at least $p-2\varepsilon p-\varepsilon p-3\ge n'$ usable
nodes while $\mathrm{disc}<\varepsilon p$, and by the bookkeeping it is admissible
conditionally on $\mathcal H_f$. By Lemma~\ref{lem:er-l1}(b) for both
explorations, outside probability $2Cp^{-8}$: if $T^{\mathrm{out}}_v,T^{\mathrm{in}}_w\ge t_0$
and no connection occurred, then $|A_f|,|A_b|\ge\delta p$. $A_f$ and $A_b$ are
disjoint, and every pair $(a,a')\in A_f\times A_b$ is unrevealed (its tail is not
in $X_f$, its head not in $X_b$), so given both histories these pairs are
independent Be$(q)$ (or forced present): $P(\text{no edge }A_f\to A_b)\le(1-q)^{\delta^2p^2-3}\le e^{-c\delta^2p/2}$.
An edge $a\to a'$ gives $v\leadsto a\to a'\leadsto w$. Union bound over $p^2$
pairs.
\end{proof}

Let $\Omega$ be the intersection of the events of Lemma~\ref{lem:er-l1}
(consequence) and Lemma~\ref{lem:er-l2}, $P(\Omega^c)\le Cp^{-5}$. Put
$\mathcal S=\{v:T^{\mathrm{out}}_v\ge t_0\}$, $\mathcal T=\{v:T^{\mathrm{in}}_v\ge t_0\}$.

\begin{corollary}[structure]\label{cor:er-struct}
On $\Omega$: $\mathcal S\times\mathcal T$ is contained in the reachability
relation; every SCC other than the one containing $\mathcal S\cap\mathcal T$ has
fewer than $t_0$ nodes; if $|\mathcal S\cap\mathcal T|\ge t_0$, this SCC equals
$\mathcal S\cap\mathcal T$.
\end{corollary}
\begin{proof}
The first claim is Lemma~\ref{lem:er-l2}. A node $u\notin\mathcal S\cap\mathcal T$
has $T^{\mathrm{out}}_u<t_0$ or $T^{\mathrm{in}}_u<t_0$, and its SCC lies in
$\mathrm{out}(u)\cap\mathrm{in}(u)$. If $u$ were in the SCC of
$\mathcal S\cap\mathcal T$, that SCC (of size $\ge t_0$) would lie in
$\mathrm{out}(u)\cap\mathrm{in}(u)$, a contradiction.
\end{proof}
(This re-proves the part of the Karp--{\L}uczak structure theorem we need; the
statement is recorded in the introduction of \citealp{dovgal2020birth}.)

\begin{lemma}[local concentration; standard coupling]\label{lem:er-local}
Fix $R$ and a function $\varphi$ of the isomorphism class of the radius-$R$
neighbourhood $B_R(e)$ of a non-loop edge $e$ (all edges incident to nodes at
undirected distance $<R$ from $e$, rooted at the oriented edge $e$), and let
$\Phi_p=p^{-1}\sum_e\varphi(B_R(e))$, the sum over non-loop edges.
(a) If $\varphi$ is bounded, $E\Phi_p\to cE_{\mathcal G}\varphi$ and
$\mathrm{Var}\,\Phi_p\to0$; the same holds for node-rooted functionals with the
node-rooted Poisson tree.
(b) If $\varphi\ge0$ and $\varphi(B_R(e))\le\mathrm{indeg}(j)$ for $e=(i,j)$, then
$\limsup_pE\Phi_p\le cE_{\mathcal G}[\min(\varphi,M)]+C/M$ for every $M$, with $C$
not depending on $p$ or $M$.
\end{lemma}
\begin{proof}
(a) First moment:
$E\sum_e\varphi(B_R(e))=\sum_{(i,j)}q\,E[\varphi(B_R((i,j);G\cup\{(i,j)\}))]$.
Explore the undirected neighbourhood of $(i,j)$ breadth first, revealing for
each explored node $w$ all pairs $(w,u)$ and $(u,w)$. On $\{|B_R|\le M\}$, each
explored node has $\mathrm{Bin}(|U|,q)$ new children and $\mathrm{Bin}(|U|,q)$
new parents among the unseen nodes, $|U|\ge p-M$, which (F1) couples with
Poisson$(c)$ up to an error $\le Mq^2p+cM/p$ per count, and the pairs between
already seen nodes are present with total probability $\le M^2q$. So the
truncated ball is within $C(M,R)/p$ in total variation of the truncated
$\mathcal G$. Tightness: each exploration step is dominated by
$\mathrm{Bin}(2p,q)$ (resp.\ Poisson$(2c)$), so $E|B_R|\le(1+2c+o(1))^R$ and
$P(|B_R|>M)\le C^R/M$ uniformly in $p$; let $p\to\infty$, then $M\to\infty$.
Second moment: $E\Phi_p^2=p^{-2}\big[\sum_e+\sum_{e,f\text{ sharing a node}}+\sum_{e,f\text{ node-disjoint}}\big]$.
The first two sums are $O(p)$ since $E\sum_vd(v)^2=O(p)$. For node-disjoint
$e,f$, explore both balls jointly in $G\cup\{e,f\}$: on $\{$both sizes $\le M\}$
they are disjoint and non-adjacent except with probability $\le C(M,R)/p$, and
then the two explorations are coupled with independent copies of $\mathcal G$.
Hence $p^{-2}\sum_{\text{disjoint}}\to(cE_{\mathcal G}\varphi)^2$ and
$\mathrm{Var}\,\Phi_p\to0$.
(b) Apply (a) to $\min(\varphi,M)$ and bound the rest by
$E\sum_{(i,j)}\mathrm{indeg}(j)\mathbf 1[\mathrm{indeg}(j)>M]=E\sum_j\mathrm{indeg}(j)^2\mathbf 1[\mathrm{indeg}(j)>M]\le p\,E[D^3]/M$,
$D\sim\mathrm{Bin}(p-1,q)$, and $E[D^3]\le c^3+3c^2+c$ for all $p$ (uniform
integrability).
\end{proof}

\textbf{Local functionals.} All of the following are functions of $B_{2K+1}(e)$
(a breadth-first search of the first $K$ nodes of a set stays within distance
$K-1$), are invariant under isomorphism, and count all qualifying nodes, so no
labelling enters.
$A^K_e=\mathbf 1[T^{\mathrm{out}}_j\ge K,\ T^{\mathrm{in}}_i\ge K]$, with
$E_{\mathcal G}A^K=\theta_K^2$ (the descendant tree of $j$ and the ancestor tree
of $i$ are independent Poisson Galton--Watson trees).
$\Psi_K(e)=\mathrm{cl}_1\vee\mathrm{cl}_2$, with
$\mathrm{cl}_1=[T^{\mathrm{in}}_i<K$ and $j$ has a parent $k\notin\{i,j\}]$ and
$\mathrm{cl}_2=[$some in-isolated $s$ has $T^{\mathrm{out}}_s<K$, $j\in\mathrm{out}(s)$, $i\notin\mathrm{out}(s)]$.
$\varphi^{\beta}_K(e)=\#\{k\notin\{i,j\}:k\to j,\ T^{\mathrm{in}}_k<K,\ T^{\mathrm{in}}_i<K,\ \mathrm{in}(k)\cap\mathrm{in}(i)\ne\emptyset\}\le\mathrm{indeg}(j)$,
with $E_{\mathcal G}\min(\varphi^\beta_K,M)=0$ (in $\mathcal G$ the ancestor trees
of $i$ and of another parent of $j$ are disjoint).

\begin{proposition}[$L$, $c>1$]\label{prop:er-N}
For $c>1$, $N/p\to cx^2$ in $L^1$ and in probability.
\end{proposition}
\begin{proof}
By (W5), $N=\sum_eZ_e$ with $Z_e=\mathbf 1[j\leadsto i]$; compare with
$\tilde N_K=\sum_eA^K_e$.
(i) $Z=1,A=0$: $j\leadsto i$ and ($T^{\mathrm{out}}_j<K$ or $T^{\mathrm{in}}_i<K$).
By (W5) these events do not depend on the presence of $e$, so the expected
number of such edges is $p(p-1)q\,P(i\in\mathrm{out}(j),T^{\mathrm{out}}_j<K)$ plus
the same with $\mathrm{in}(i)$, and by exchangeability of the $p-1$ nodes other
than $j$, $P(i\in\mathrm{out}(j),T^{\mathrm{out}}_j<K)=E[(T^{\mathrm{out}}_j-1)\mathbf 1(T^{\mathrm{out}}_j<K)]/(p-1)\le K/(p-1)$.
Total $\le2cK$.
(ii) $Z=0,A=1$: on $\Omega$, $T^{\mathrm{out}}_j\ge t_0$ and $T^{\mathrm{in}}_i\ge t_0$
would give $j\leadsto i$, so $T^{\mathrm{out}}_j$ or $T^{\mathrm{in}}_i$ lies in
$[K,t_0)\subseteq[K,\varepsilon p)$. The expected number of such edges is at
most $\sum_{(i,j)}q\,P_{G^{\{e\}}}(\cdot)\le cp\cdot2Ce^{-\kappa K}$ by
Lemma~\ref{lem:er-l1}(a) in $G^{\{e\}}$; on $\Omega^c$ the count is at most $m$
and $E[m\mathbf 1_{\Omega^c}]\le(Em^2)^{1/2}P(\Omega^c)^{1/2}=O(p\cdot p^{-5/2})$.
So $E|N-\tilde N_K|\le2cK+2cCpe^{-\kappa K}+o(1)$. By
Lemma~\ref{lem:er-local}(a), $\tilde N_K/p\to c\theta_K^2$ in $L^2$; hence
$\limsup_pE|N/p-cx^2|\le2cCe^{-\kappa K}+c|\theta_K^2-x^2|\to0$ as $K\to\infty$.
\end{proof}
The same argument for the node functional $\mathbf 1[T^{\mathrm{out}}_v\ge K,T^{\mathrm{in}}_v\ge K]$
gives $|\mathcal S\cap\mathcal T|/p\to x^2$, so by \cref{cor:er-struct} the giant
SCC has $(x^2+o_P(1))p$ nodes and all other SCCs fewer than $C_3\log p$, w.h.p.

\begin{proposition}[the tree computation]\label{prop:er-tree}
In $\mathcal G$, with $\Psi_\infty$ defined by $T^{\mathrm{in}}_i<\infty$ and
$T^{\mathrm{out}}_s<\infty$, $E_{\mathcal G}\Psi_K\uparrow P_{\mathcal G}(\Psi_\infty)=F_\uparrow$
as $K\to\infty$.
\end{proposition}
\begin{proof}
Both clauses increase in $K$ to their $K=\infty$ versions (monotone convergence).
Let $A(i)$ be the event that the ancestor tree of $i$ is infinite and $D(j)$ that
the descendant tree of $j$ is infinite (probability $x$ each, independent). The
other parents of $j$ are Poisson$(c)$ many.
(1) Not $A(i)$: $\Psi_\infty$ holds iff $j$ has another parent. Clause 1 needs
exactly that; clause 2 needs it too, since a path from $s$ to $j$ ends with an
edge $k\to j$ and $k=i$ would put $i$ into $\mathrm{out}(s)$. Probability
$(1-x)(1-e^{-c})$.
(2) $A(i)$: clause 1 fails. Clause 2 needs (a) not $D(j)$, since
$\mathrm{out}(s)$ contains the descendant tree of $j$ (probability $1-x$), and
(b) another parent $k$ of $j$ and a chain $k=u_0\leftarrow u_1\leftarrow\dots\leftarrow u_r=s$
ending at an in-isolated $s$, in which every $u_t$ has only finite descendant
trees below its children other than $u_{t-1}$ ($u_{-1}=j$); such an $s$ does not
reach $i$, because paths into $i$ stay in the ancestor tree of $i$. The
probability that all other children of a node have finite descendant trees is
$E[(1-x)^{\mathrm{Poi}(c)}]=e^{-cx}=1-x$. Let $Q(v)$ be the event that $v$ is
in-isolated or some parent $u$ of $v$ has all other children finite and $Q(u)$
holds, and $d^*=P(Q)$. Parents are Poisson$(c)$ many and each is \emph{good}
(other children finite and $Q$) independently with probability $(1-x)d^*$, so by
Poisson thinning $d^*=e^{-c}+1-e^{-c(1-x)d^*}=f(d^*)$. $Q$ is the increasing
union of its depth-$n$ truncations $Q_n$, with $P(Q_0)=e^{-c}=f(0)$ and
$P(Q_{n+1})=f(P(Q_n))$; as $f$ is increasing and continuous, $P(Q_n)=f^{n+1}(0)$
increases to the smallest fixed point $d$ of $f$ (if $d'$ is a fixed point,
$f^n(0)\le f^n(d')=d'$). So $d^*=d$. The number of good other parents of $j$ is
Poisson$(c(1-x)d)$, so (b) has probability $1-e^{-c(1-x)d}$, independently of
$A(i)$ and $D(j)$. Adding (1) and (2) gives $F_\uparrow$.
\end{proof}
(An independent Monte Carlo of the event itself, not of the closed form, gives
$0.38756\pm0.00077$ against $0.38773$ at $c=1.5$ and $0.19049\pm0.00062$ against
$0.18962$ at $c=2$.)

\begin{proposition}[reduction of $\Eu$ to $\Psi_K$, $c>1$]\label{prop:er-red}
$E\sum_e|\mathbf 1[e\in\Eu]-\Psi_K(e)|\le(\delta_1(K)+o(1))p$ with
$\delta_1(K)=C(c^2+K^2)e^{-\kappa K}+o_K(1)$, where $o_K(1)\to0$ as $K\to\infty$
and $C$ does not depend on $K$ or $p$.
\end{proposition}
\begin{proof}
Let $e=(i,j)$ be a non-loop edge. \emph{Clause 2 never errs}: an in-isolated $s$
is a source node, so $j\in\mathrm{out}(s)$, $i\notin\mathrm{out}(s)$ give
$e\in\Eu$. \emph{Case I: $T^{\mathrm{in}}_i<K$.} (I-a) If $e\in\Eu$, (W1) gives
another parent of $j$, so clause 1 holds: no error. (I-b) Clause 1 holds and
$e\notin\Eu$. Then
$\mathbf 1[\text{I-b}]\le\sum_{k\notin\{i,j\}}\mathbf 1[e,(k,j)\text{ present},\ T^{\mathrm{in}}_i<K,\ e\notin\Eu]$
(a union over \emph{all} other parents; no parent is singled out). Fix the
triple $(i,j,k)$ and force $F=\{(i,j),(k,j)\}$. By (W3), $e\notin\Eu$ forces
$\mathrm{in}(k)\cap\mathrm{in}(i)\ne\emptyset$, and by (W2) every source node
reaching $k$ reaches $i$. Split on $T^{\mathrm{in}}_k$:
($\alpha_1$) $K\le T^{\mathrm{in}}_k<\varepsilon p$: probability $\le Ce^{-\kappa K}$
(Lemma~\ref{lem:er-l1}(a), backward, in $G^F$).
($\alpha_2$) $T^{\mathrm{in}}_k\ge\varepsilon p$: the backward exploration of
$\mathrm{in}(k)$ explores at least $\varepsilon p/2$ nodes. When a node $y\ne j$ is
explored, the pairs $(u,y)$ with $u$ unseen are revealed (all absent with
conditional probability $\ge(1-q)^p$), and the pairs $(u,y)$ with $u$ already
seen are never revealed by this exploration and are independent of it and of
each other (distinct heads); so, given the history, $y$ is in-isolated with
probability $\ge(1-q)^p\ge e^{-c}/2$, and the number of in-isolated nodes among
the first $\varepsilon p/2$ explored nodes dominates
$\mathrm{Bin}(\varepsilon p/2-1,e^{-c}/2)$. By (F2) at least $K$ of them are
in-isolated except with probability $e^{-\Omega(p)}$. Each in-isolated node of
$\mathrm{in}(k)$ is a source node reaching $k$, hence reaches $i$, so lies in
$\mathrm{in}(i)$; but $|\mathrm{in}(i)|<K$. So ($\alpha_2$) has probability
$e^{-\Omega(p)}$.
($\beta$) $T^{\mathrm{in}}_k<K$ and $\mathrm{in}(k)\cap\mathrm{in}(i)\ne\emptyset$.
Summed over $k$ this is $\varphi^\beta_K(e)$.
The number of triples is $p^3$ and each has both pairs present with probability
$q^2$, so cases ($\alpha_1$), ($\alpha_2$) contribute at most
$c^2p(Ce^{-\kappa K}+e^{-\Omega(p)})$; case ($\beta$) contributes
$E\sum_e\varphi^\beta_K(e)\le(C/M+o(1))p$ for every $M$ by
Lemma~\ref{lem:er-local}(b), i.e.\ $o(p)$.
\emph{Case II: $T^{\mathrm{in}}_i\ge K$.} Clause 1 fails.
(II-0) $K\le T^{\mathrm{in}}_i<\varepsilon p$: at most $cpCe^{-\kappa K}$ such
edges in expectation (Lemma~\ref{lem:er-l1}(a) in $G^{\{e\}}$).
(II-1) $T^{\mathrm{in}}_i\ge\varepsilon p$, so $i\in\mathcal T$. Work on
$\Omega'=\Omega\cap\{\max_v\mathrm{indeg}(v)\le C\log p\}$; $P(\Omega'^c)\le Cp^{-2}$,
and the edges counted on $\Omega'^c$ contribute at most
$E[m\mathbf 1_{\Omega'^c}]\le(Em^2)^{1/2}P(\Omega'^c)^{1/2}=O(1)$. An error here is
an edge $e\in\Eu$ with $\Psi_K(e)=0$. Take a witness source node $z$ of (W1),
in a source SCC $\sigma$. On $\Omega$, $z\in\mathcal S$ would give $z\leadsto i$
(Lemma~\ref{lem:er-l2}), so $T^{\mathrm{out}}_z<t_0$.
(A) $\sigma=\{z\}$ (i.e.\ $z$ in-isolated) and $T^{\mathrm{out}}_z<K$: clause 2
holds with $s=z$; no error.
(B1) $z$ in-isolated and $K\le T^{\mathrm{out}}_z<t_0$. The number of such edges
is at most $\sum_z\mathbf 1[K\le T^{\mathrm{out}}_z<t_0]\sum_{j\in\mathrm{out}(z)}\mathrm{indeg}(j)$.
Given $\mathrm{out}(z)=O$ with $|O|=r$ (a completed forward exploration), the
pairs with tail outside $O$ and head in $O$ are unrevealed, so the expected
number of edges into $O$ is at most $r^2+rpq=r^2+cr$. With
$P(T^{\mathrm{out}}_z=r)\le Ce^{-\kappa r}$ (Lemma~\ref{lem:er-l1}(a)), the
expectation is at most $p\sum_{r\ge K}Ce^{-\kappa r}(r^2+cr)\le C'pK^2e^{-\kappa K}$.
(B2) $|\sigma|\ge2$. Then $|\sigma|<t_0$ (as $\sigma\subseteq\mathrm{out}(z)$),
every node of $\sigma$ has out-set $\mathrm{out}(z)$ of size $<t_0$, and the
number of edges with head in $\mathrm{out}(z)$ is at most $t_0\cdot C\log p$ on
$\Omega'$. Let $N_\sigma$ be the number of source SCCs with $2\le|\sigma|<t_0$.
For a set $R$ of size $r$, ``$R$ strongly connected'' and ``no edge enters $R$''
depend on disjoint sets of pairs; the first requires every node of $R$ to have
a parent in $R$ (probability $\le(rq)^r$, distinct heads), the second has
probability $(1-q)^{r(p-r)}$. Hence
$EN_\sigma\le\sum_{r=2}^{t_0}\binom pr(rq)^r(1-q)^{r(p-r)}\le\sum_{r\ge2}\big(ce^{1-c}e^{cr/p}\big)^r=O(1)$,
using $\binom pr\le p^r/r!$, $r^r/r!\le e^r$, $r\le t_0=O(\log p)$ and
$ce^{1-c}<1$ for $c\ne1$. So the (B2) errors have expectation
$O(1)\cdot t_0C\log p=O(\log^2p)=o(p)$: an expectation on the event $\Omega'$ plus
the Cauchy--Schwarz bound for $\Omega'^c$ above.
Adding all cases gives the claim.
\end{proof}

\begin{theorem}[$S$, $c>1$]\label{thm:er-S}
For $c>1$, $|\Eu|/p\to cF_\uparrow$, and therefore
$S=\tfrac12+|\Eu|/(2(p+m))\to\tfrac12+cF_\uparrow/(2(1+c))$, in $L^1$ and in
probability.
\end{theorem}
\begin{proof}
$E\big||\Eu|/p-cF_\uparrow\big|\le E\sum_e|\mathbf 1[e\in\Eu]-\Psi_K(e)|/p+E|\Phi^{\Psi_K}_p-cE_{\mathcal G}\Psi_K|+c|E_{\mathcal G}\Psi_K-F_\uparrow|$;
the terms are bounded by Prop.~\ref{prop:er-red}, Lemma~\ref{lem:er-local}(a) and
Prop.~\ref{prop:er-tree}; let $p\to\infty$, then $K\to\infty$. $m/p\to c$ in
$L^2$ ($m\sim\mathrm{Bin}(p(p-1),q)$), and $S$ and $L$ are bounded continuous
functions of $(|\Eu|/p,m/p)$ and $(N/p,m/p)$ on $[0,\infty)^2$.
\end{proof}

\begin{proposition}[$c\le1$]\label{prop:er-sub}
(a) For $c<1$, $EN\le c^2/(1-c)$. (b) For $c\le1$, $N/p\to0$ in $L^1$, so
$L\to\tfrac12+\frac{c}{2(1+c)}$ ($L(1)=3/4$). (c) For $c\le1$,
$|\Eu|/p\to c(1-e^{-c})$ in $L^1$, so $S\to\tfrac12+\frac{c(1-e^{-c})}{2(1+c)}$
($S(1)=0.65803$); this is $F_\uparrow$ at $x=0$.
\end{proposition}
\begin{proof}
(a) A non-loop edge inside an SCC lies on a simple directed cycle, and the
expected number of pairs (simple cycle of length $k$, edge on it) is
$k\cdot\frac{(p)_k}{k}q^k\le c^k$; sum over $k\ge2$.
(b) $c<1$ from (a). $c=1$: $D(p,1/p)$ is a subgraph of $D(p,c'/p)$ for every
$c'>1$ under the monotone coupling, and an edge inside an SCC of the subgraph is
inside an SCC of the larger graph, so $\limsup_pEN/p\le c'x(c')^2\to0$ as
$c'\downarrow1$ (Prop.~\ref{prop:er-N}).
(c) Let $\Psi^1_K=\mathrm{cl}_1$. Case I ($T^{\mathrm{in}}_i<K$): (I-a) as before,
no error; (I-b) with the same union over other parents $k$ and forced
$F=\{(i,j),(k,j)\}$, (W3) forces $\mathrm{in}(k)\cap\mathrm{in}(i)\ne\emptyset$,
so either $T^{\mathrm{in}}_k\ge K$ or the event of $\varphi^\beta_K$ holds. The
expected number of triples with $T^{\mathrm{in}}_k\ge K$ is
$E\sum_e\#\{k\notin\{i,j\}:k\to j,\ T^{\mathrm{in}}_k\ge K\}$, a local functional
dominated by $\mathrm{indeg}(j)$, so by Lemma~\ref{lem:er-local}(b) it is at most
$(c\cdot c\,\theta_K+C/M+o(1))p$; the $\varphi^\beta_K$ part is $o(p)$ as before.
Case II ($T^{\mathrm{in}}_i\ge K$): at most $E\sum_e\mathbf 1[T^{\mathrm{in}}_i\ge K]=(c\theta_K+o(1))p$
(Lemma~\ref{lem:er-local}(a)). Hence
$\limsup_pE\sum_e|\mathbf 1[e\in\Eu]-\Psi^1_K(e)|/p\le c^2\theta_K+c\theta_K+C/M$, which
tends to $0$ as $K,M\to\infty$ because $\theta_K\to x(c)=0$ for $c\le1$ (F3).
Finally $E_{\mathcal G}\Psi^1_K=(1-\theta_K)(1-e^{-c})\to1-e^{-c}$.
\end{proof}

\textbf{What is not proved.} Rates of convergence; the critical window
$c=1+\lambda p^{-1/3}$ (\L{}uczak and Seierstad, recorded in
\citealp{dovgal2020birth}; not used); almost sure convergence; the regime of
fixed $T$ with $p\to\infty$. The proof does not use the structure theorems of
Karp and of \L{}uczak (not opened by us); it re-proves the part it needs.

\textbf{Numerical checks of the intermediate claims} (fixed seeds, exact
computation; supplement \texttt{x1proof/verify}). (W4) $\Eu\subseteq$ edges into
confluences: 0 violations in 48 graphs. Dichotomy (Lemma~\ref{lem:er-l1}): at
$c=2$, $p=10^5$, the mass of $K\le T<0.02p$ is $0.0097$, $0.0008$, $0$ for
$K=5,10,20$; at $c=1$ it is $0.37$, $0.25$, $0.17$ (critical, no dichotomy, as
expected). Structure (\cref{cor:er-struct}): for $c\ge1.5$ the second-largest
SCC never had more than 2 nodes; giant fraction $0.63439$ ($p=10^6$, $c=2$)
against $x^2=0.63491$. Prop.~\ref{prop:er-red} edge by edge: at $c=2$,
$p=2\cdot10^5$, $|\Eu|=76{,}220$ and $\sum_e\Psi_{40}(e)=76{,}220$ with no edge
in either difference; clause 2 never erred in any run. Concentration: the
standard deviation of $N/p$ times $\sqrt p$ is $2.88$--$3.07$ for
$p=5\cdot10^3$ to $3.2\cdot10^5$ ($c=2$), as Lemma~\ref{lem:er-local} requires.

%=============================================================================
\section{Accuracy Interpretation of $\srho$}
\label{app:x3}

An edge $\{x,y\}$ of a DAG is shown in a uniformly random order together with
$(\rho_{\mathrm{first}},\rho_{\mathrm{second}})$. A \emph{pairwise count rule} is
$f(a,b)\in[0,1]$, the probability of declaring ``first is the tail''; the
\emph{sort rule} is $f_{\mathrm{sort}}=1,\tfrac12,0$ for $a<,=,>b$. For the true
edge $x\to y$ the accuracy is $\tfrac12f(\rho_x,\rho_y)+\tfrac12(1-f(\rho_y,\rho_x))$.

\begin{proposition}[elementary]\label{prop:x3}
On every DAG, (i) the accuracy of $f_{\mathrm{sort}}$ is exactly $\srho$;
(ii) every pairwise count rule has accuracy at most $\srho$; (iii)
$f_{\mathrm{sort}}$ never errs on an edge with $\rho_x\ne\rho_y$.
\lean{\lname{accRule_sortRule}, \lname{accRule_sortRule_rho}, \lname{accRule_le_sOf}, \lname{accRule_le_sRho}, \lname{edgeAcc_sortRule_strict}, \lname{rho_not_gt_of_edge}.}
\end{proposition}
\begin{proof}
By Reisach et al.'s Theorem~1, $\rho_x\le\rho_y$ on an edge $x\to y$. If
$\rho_x<\rho_y$ the accuracy is at most $1$, attained by the sort rule; if
$\rho_x=\rho_y=a$ it is $\tfrac12f(a,a)+\tfrac12(1-f(a,a))=\tfrac12$ for every
$f$, because the random order forces the tie decision to $1/2$. These are the edge
scores. (i) restates the definition of $\srho$; (ii) is a three-line consequence of
their Theorem~1.
\end{proof}
Combined with \cref{thm:stat}, the best accuracy of any pairwise count rule on the
stationary unrolled DAG is $\tfrac12+|\Eu|/(2|\Es|)$, at most $0.665$ for ER
summary graphs and $0.75$ for river-like in-forests ($\mu\le1$). The restriction
to counts is necessary (exhaustive enumeration over labelled DAGs, not proved):

\begin{center}\small
\begin{tabular}{lcccc}
\toprule
nodes $p$ & 2 & 3 & 4 & 5 \\
\midrule
Bayes rule on the two counts (= edge-pooled $\srho$) & 0.5 & 0.5625 & 0.592262 & 0.605547 \\
Bayes rule on counts and the two node degrees & 0.5 & 0.625 & 0.660714 & 0.669238 \\
\bottomrule
\end{tabular}
\end{center}
For the skeleton $a-c-b$ the collider has counts $(2,2,3)$ (all edges strict)
while the chain and the fork have all counts equal to $3$; among the tied graphs the
centre $c$ is the tail with probability $2/3$ under a uniform prior, so a
degree-aware rule beats the coin. So $\srho$ is the Bayes accuracy of pairwise
count rules only. This section is interpretive, not a new theorem.
